% !TeX root = ../../Book3.tex
% 纲要标题：## 第20章　Hausdorff 测度、可整流性与面积公式
% 纲要位置：几何学计划纲要.md，第 2922 行。
% 正文状态：初稿；尚未正式审核。

% 纲要内容开始
% 
% > 本章目标与证明范围：先建立可测几何的真正局部工具，再引入弱曲面。Rademacher、面积和余面积各有独立证明步骤，近似切平面不是光滑切平面的名义推广。
% 
% 纲要内容结束

\chapter{Hausdorff 测度、可整流性与面积公式}
\label{g3:ch:20}
\BookChapterNumber{g3:ch:20}

\begin{chapterabstract}
	光滑曲面的面积来自切向量；当曲面出现折痕、重叠或细碎的可测部分时，需要先建立不依赖光滑坐标的测度。本章从 Hausdorff 外测度与覆盖引理出发，证明它在 Euclid 平面及光滑子流形上与通常面积一致。绝对连续函数的积分表示引出 Rademacher 定理，可数一致线性化随后把一般 Lipschitz 映射分解为可计算的片。面积公式记录像点的重数，余面积公式沿纤维积分，两者分别完成证明。最后用参数片的放大极限建立可整流集的近似切平面、密度和重数，并构造具有正有限长度的纯不可整流集。
\end{chapterabstract}

\begin{learninggoals}
	\begin{itemize}
		\item 解释 Hausdorff 测度的归一化，使用五倍球与 Vitali 覆盖，证明密度定理及光滑面积的一致性。
		\item 从一维绝对连续性证明 Lipschitz 映射几乎处处可微，区分方向导数、全微分与近似微分。
		\item 把逐点微分整理为可数一致线性化片，说明满秩如何产生双 Lipschitz 参数。
		\item 证明带权面积公式，核对重数可测性，并区分几何面积、参数面积与推前质量。
		\item 证明余面积公式，掌握零测集、秩亏集与满秩片各自的切片处理方法。
		\item 建立可整流集的近似切平面与密度结论，识别整数重数和纯不可整流部分。
	\end{itemize}
\end{learninggoals}

把一张纸沿直线折起，面积仍然可以计算，但折痕上的切平面已经不再唯一。若把一条曲线描画两遍，参数长度加倍，像集本身却没有变大。若再把许多小片拼在一起，集合甚至可能在每个尺度上都向不同方向散开。这些现象提示我们分别回答三个问题：怎样直接测量一个集合，怎样由参数计算其测度，以及什么条件能使几乎每一点仍具有切向线性模型。

\ProseParagraphBreak

对光滑单射参数 \(F:\Omega\subset\mathbb R^m\to\mathbb R^n\)，已有的面积表达式是
\[
	\int_\Omega\sqrt{\det\bigl((\mathrm DF)^{\mathsf T}\mathrm DF\bigr)}\,\mathrm dx.
\]
本章将逐步去掉光滑性和单射性。Lipschitz 条件保证微分几乎处处存在；当像点对应多个原像时，把这些原像数加起来，就得到参数面积。源维数大于目标维数时，原像通常不再是离散点，而是低维切片，于是计数要换成纤维上的 Hausdorff 测度。这两种积分公式需要不同的秩亏估计，不能只凭光滑换元公式直接推出。

\ProseParagraphBreak

证明顺序从覆盖出发：先获得测度微分工具，再证明可微性，继而在近似线性的片上计算面积与切片。近似切平面的存在放在这些公式之后，因为它的证明要把参数中的线性近似转化为像上的测度收敛。全章沿用\chapref{g3:ch:11}的基本 Lebesgue 积分、测度正则性、Fubini 定理、收敛定理和绝对连续测度的积分密度表示；分布只用于识别方向导数的线性关系，并会重述所需唯一性论证。线性代数使用正交分解与对称矩阵对角化。光滑面积与\chapref{g3:ch:19}衔接，但本章的主要测度与微分结果各有证明。以下五节均属主线。

% !TeX root = ../../Book3.tex
% 纲要标题：### 20.1 测度及覆盖方法
% 纲要位置：几何学计划纲要.md，第 2926 行。
% 正文状态：初稿；尚未正式审核。

% 纲要内容开始
% 
% * Hausdorff 外测度与归一化
% * Vitali 型覆盖引理
% * 测度微分与密度定理
% * Hausdorff 测度与子流形体积
% 
% 纲要内容结束

\section{测度及覆盖方法}
\label{g3:sec:2001}

一条曲线在平面中的二维 Lebesgue 测度为零，却可以有非零长度。要在同一环境中同时测量曲线、曲面和不光滑的极限集合，就需要由集合自身的小尺度覆盖来定义不同维数的测度。本节先建立这一测量方法，再证明覆盖、密度与局部积分之间的联系。所用基本测度知识与\chapref{g3:ch:11}相同，包括积分收敛与换序、有限测度的正则性以及绝对连续测度的积分密度表示；覆盖和微分结论在这里证明。

\subsection{Hausdorff 外测度与归一化}
\label{g3:subsec:2001:01}
用许多直径很小的集合覆盖一个对象，将每个覆盖集的直径取 \(s\) 次方再求和，就得到一个 \(s\) 维大小的候选值。细化覆盖会排除用一只大球掩盖局部结构的做法。整数维数中的常数则应使平面单位球具有通常的体积。

\begin{definition}\label{g3:def:hausdorff-measure}
	设 \((X,d)\) 为度量空间，\(s\ge0\)，记
	\(\omega_s=\pi^{s/2}/\Gamma(1+s/2)\)，其中 \(\Gamma\) 是 Gamma 函数；当 \(s=m\) 为正整数时，\(\omega_m\) 是 \(\mathbb R^m\) 单位球的 Lebesgue 体积。对 \(E\subset X\)、\(\delta>0\)，定义
	\[
		\mathcal H^s_\delta(E)=
		\inf\left\{\sum_j\omega_s\left(\frac{\operatorname{diam}U_j}{2}\right)^s:
		E\subset\bigcup_jU_j,\ \operatorname{diam}U_j<\delta\right\},
	\]
	下确界遍历可数覆盖，空集的代价为零；\(s=0\) 时每个非空覆盖集的代价为一。相应的\term[Hausdorff cedu]{Hausdorff 测度}[Hausdorff measure]为
	\begin{equation}\label{g3:eq:hausdorff-measure-definition}
		\mathcal H^s(E)=\lim_{\delta\downarrow0}\mathcal H^s_\delta(E).
	\end{equation}
	这个极限按单调增加取得，允许为无穷大。在未说明集合可测以前，\(\mathcal H^s\) 表示外测度。
\end{definition}
\symidx{\(\mathcal H^s\)：采用单位球体积归一化的 \(s\) 维 Hausdorff 测度}

零维定义给出计数测度：有限个不同点在足够小的尺度下必须分别覆盖；无限集包含任意多个不同点，所以其零维测度为无穷大。正维定义则使单点为零测集。\(\mathcal H^s\) 总是相对于所给度量定义，改变度量可能改变数值及维数。

\begin{proposition}\label{g3:prop:hausdorff-borel-measure}
	对任意度量空间 \((X,d)\) 及 \(s\ge0\)，式\eqref{g3:eq:hausdorff-measure-definition}定义一个外测度；若 \(A,B\subset X\) 满足 \(\operatorname{dist}(A,B)>0\)，则
	\[
		\mathcal H^s(A\cup B)=\mathcal H^s(A)+\mathcal H^s(B).
	\]
	所有 Borel 集均可测，任意集合都包含于同外测度的 Borel 集。因此将 \(\mathcal H^s\) 限制在其可测集上，得到一个完备测度。
\end{proposition}
\begin{proof}
	空集和单调性直接来自覆盖。把每个 \(E_i\) 的近似最优覆盖合并，得到
	\(\mathcal H^s_\delta(\bigcup_iE_i)\le\sum_i\mathcal H^s_\delta(E_i)\)；右侧再以 \(\sum_i\mathcal H^s(E_i)\) 控制，令 \(\delta\downarrow0\) 即得可数次可加性。当 \(\delta<\operatorname{dist}(A,B)\) 时，一个覆盖集不能同时遇到两者，把覆盖按相交对象分开便得反向不等式，进而得到距离为正的可加性。

	说明这一距离性质为何保证 Borel 可测。记外测度为 \(\mu\)，给定闭集 \(F\) 和满足 \(\mu(T)<\infty\) 的集合 \(T\)，令
	\[
		T_j=\{x\in T:d(x,F)\ge1/j\},\qquad
		C_j=\{x\in T:1/(j+1)\le d(x,F)<1/j\}.
	\]
	有 \(\mu(T)\ge\mu(T\cap F)+\mu(T_j)\)。有限个奇数编号的 \(C_j\) 两两距离为正，偶数编号同样如此；故每组的测度和均至多为 \(\mu(T)\)，从而 \(\sum_j\mu(C_j)<\infty\)。又
	\(T\setminus F\subset T_j\cup\bigcup_{i\ge j}C_i\)，所以
	\(\mu(T\setminus F)\le\lim_j\mu(T_j)\)。这给出
	\(\mu(T)\ge\mu(T\cap F)+\mu(T\setminus F)\)；另一方向来自次可加性。若 \(\mu(T)=\infty\)，这一不等式自动成立。由外测度的可测性判据，闭集可测，因而其生成的 Borel 集均可测。

	最后设 \(\mathcal H^s(E)<\infty\)。在尺度 \(1/j\) 选择覆盖，使总代价至多为 \(\mathcal H^s(E)+1/j\)，并把各覆盖集取闭包。令 \(B\) 为这些可数闭集并的可数交，则 \(B\) 是包含 \(E\) 的 Borel 集，而且
	\(\mathcal H^s_{1/j}(B)\le\mathcal H^s(E)+1/j\)。令 \(j\to\infty\) 得到相等测度。若外测度为无穷大，可直接取 \(B=X\)。外测度的全部可测集组成完备测度空间，这是同一可测性判据的零测子集情形。
\end{proof}

\begin{proposition}\label{g3:prop:hausdorff-lipschitz}
	设 \((X,d_X)\)、\((Y,d_Y)\) 为度量空间，\(E\subset X\)，\(f:E\to Y\) 满足
	\(d_Y(f(x),f(y))\le Ld_X(x,y)\)，其中 \(L>0\)。则对每个 \(s\ge0\) 及 \(A\subset E\)，有
	\[
		\mathcal H^s(f(A))\le L^s\mathcal H^s(A).
	\]
	这里允许不可测集合，两边均按外测度理解。特别地，等距映射保持测度；在 Euclid 空间中，\(\mathcal H^s(a+\lambda A)=\lambda^s\mathcal H^s(A)\)，\(\lambda>0\)。
\end{proposition}
\begin{proof}
	先将 \(A\) 的覆盖集与 \(E\) 相交，再取 \(f\) 的像；其直径至多乘以 \(L\)。因此
	\(\mathcal H^s_{L\delta}(f(A))\le L^s\mathcal H^s_\delta(A)\)，令 \(\delta\downarrow0\) 即得结论。等距与相似变换再对逆映射应用同一不等式。
\end{proof}

若 \(t>s\)，覆盖中的每一项满足
\[
	\omega_t(\operatorname{diam}U/2)^t
	\le\frac{\omega_t}{\omega_s}(\delta/2)^{t-s}
	\omega_s(\operatorname{diam}U/2)^s.
\]
所以 \(\mathcal H^s(E)<\infty\) 推出 \(\mathcal H^t(E)=0\)。由此定义非空集的\term[Hausdorff weishu]{Hausdorff 维数}[Hausdorff dimension]
\(\dim_{\mathrm H}E=\inf\{s\ge0:\mathcal H^s(E)=0\}\)，空集的维数在本章记为零。这个临界指数不依赖正归一化常数，但整数维的面积公式依赖它。

\subsection{Vitali 型覆盖引理}
\label{g3:subsec:2001:02}
覆盖中的球可能严重重叠。先从中选出互不相交的一族，再适度放大，就能用可加的体积和控制原覆盖。Vitali 型覆盖方法\footnote{Giuseppe Vitali（1875--1932），意大利数学家，研究实变函数、绝对连续性、覆盖定理与函数列收敛。}把这一选择用于任意小尺度的球。

\begin{lemma}\label{g3:lem:five-ball-covering}
	设 \((X,d)\) 为可分度量空间，\(\mathcal B\) 为一族正半径闭球，其半径有共同上界。存在至多可数个两两不交的球 \(B_i=\overline B_{r_i}(x_i)\in\mathcal B\)，使每个 \(B=\overline B_r(x)\in\mathcal B\) 都与某个 \(B_i\) 相交且 \(r_i\ge r/2\)。于是
	\[
		\bigcup_{B\in\mathcal B}B\subset\bigcup_i\overline B_{5r_i}(x_i).
	\]
\end{lemma}
\begin{proof}
	取半径上界 \(R>0\)，按 \(2^{-j}R<r\le2^{1-j}R\)、\(j=1,2,\ldots\) 分层。第一层选一个极大的不交子族；第 \(j\) 层在不与已选球相交的球中再选极大不交子族。极大族存在，因为按包含排序的不交子族的一条链，其并仍不交，故可用极大原理。每个不交正半径球的内部含有可数稠密集中的不同点，所以所选族至多可数。

	一个未选的第 \(j\) 层球必与本层或前层某个已选球相交，否则可加入该层的极大族。该已选球的半径大于原半径的一半。若 \(z\in B\)，相交与三角不等式给出
	\(d(z,x_i)\le2r+r_i\le5r_i\)，得到覆盖。
\end{proof}

\begin{theorem}\label{g3:thm:vitali-lebesgue-covering}
	设 \(E\subset\mathbb R^m\)、\(m\ge1\)，并给定一个包含 \(E\) 的有限 Lebesgue 测度开集 \(O\)。若闭球族 \(\mathcal B\) 包含于 \(O\)，且对每个 \(x\in E\) 和 \(\varepsilon>0\)，存在 \(B\in\mathcal B\) 满足 \(x\in B\)、\(\operatorname{diam}B<\varepsilon\)，则可选两两不交的 \(B_i\in\mathcal B\)，使
	\[
		\mathcal L^m\left(E\setminus\bigcup_iB_i\right)=0.
	\]
\end{theorem}
\begin{proof}
	可先只保留半径小于一的球，再作前一引理的选择。因这些球不交且包含于 \(O\)，有 \(\sum_i\mathcal L^m(B_i)\le\mathcal L^m(O)<\infty\)。若 \(x\in E\setminus\bigcup_iB_i\)，它到任意有限个已选闭球的并的距离为正。细覆盖条件允许取包含 \(x\)、避开前 \(N\) 个球的足够小球。前一引理的相交性质使它遇到某个编号大于 \(N\) 的已选球，且 \(x\) 落在后者的五倍球中。因此
	\[
		\mathcal L^m\left(E\setminus\bigcup_iB_i\right)
		\le5^m\sum_{i>N}\mathcal L^m(B_i)\longrightarrow0.
	\]
	这也证明结论对未预先假定可测的 \(E\) 仍按外测度成立。
\end{proof}

下面确定 Hausdorff 定义中的归一化。为得到精确常数，先把同一直径下的体积与球比较。这个短几何论证也说明为什么覆盖代价采用直径的一半。

\begin{lemma}\label{g3:lem:euclidean-isodiametric}
	对 \(A\subset\mathbb R^m\)、\(m\ge1\)，有
	\[
		\mathcal L^m(A)\le\omega_m\left(\frac{\operatorname{diam}A}{2}\right)^m,
	\]
	左侧按外测度理解。
\end{lemma}
\begin{proof}
	无界集的情形直接成立；有界集取闭包以后直径不变，故只需处理紧集 \(K\)。沿一个坐标方向，把每条平行直线上与 \(K\) 的交替换为以坐标超平面为中心、长度相同的闭区间；原交非空而长度为零时保留区间的中心点。这个操作称为 Steiner 对称化。\footnote{Jakob Steiner（1796--1863），瑞士数学家，研究射影几何、几何构造与几何极值问题。}

	先核对三个性质。纤维长度随基点上半连续：若某纤维包含于长度只多 \(\varepsilon\) 的开集，则紧性保证邻近纤维也包含于该开集，否则收敛子列将给出原纤维外的极限点。因此新集合仍紧。Fubini 定理说明它与 \(K\) 体积相同。

	对两条非空纤维，设原纤维的最小、最大坐标分别为 \(a_-,a_+\) 与 \(b_-,b_+\)，长度为 \(\ell_a,\ell_b\)。有
	\[
		\max\{|a_+-b_-|,|b_+-a_-|\}
		\ge\frac{(a_+-a_-)+(b_+-b_-)}2
		\ge\frac{\ell_a+\ell_b}2.
	\]
	新纤维中任意两点的竖直坐标差至多为最后一项，水平坐标差又未变，所以新集合直径不增加。最后，沿一个坐标方向的对称化保留其他坐标反射已有的不变性，因为反射对应的两条纤维长度相同。

	依次沿全部 \(m\) 个坐标方向处理，所得紧集 \(\widetilde K\) 体积不变、直径不增加，并关于原点中心对称。若 \(x\in\widetilde K\)，则 \(-x\in\widetilde K\)，故 \(2|x|\le\operatorname{diam}K\)。因此 \(\widetilde K\) 包含于半径 \(\operatorname{diam}K/2\) 的球，体积比较给出结论。
\end{proof}

\begin{theorem}\label{g3:thm:hausdorff-lebesgue}
	对 \(\mathbb R^m\) 上的 Euclid 度量、\(m\ge1\)，标准归一化满足
	\[
		\mathcal H^m(A)=\mathcal L^m(A)\qquad(A\subset\mathbb R^m).
	\]
	因而任意 \(m\) 维仿射平面上的 \(\mathcal H^m\) 就是该平面的通常体积。
\end{theorem}
\begin{proof}
	等直径不等式对任意小直径覆盖给出
	\(\mathcal L^m(A)\le\sum_j\omega_m(\operatorname{diam}U_j/2)^m\)；取下确界与极限得到 \(\mathcal L^m\le\mathcal H^m\)。

	反向先作一个粗估计。将 Lebesgue 外测度的立方体覆盖继续细分到直径小于 \(\delta\)，每个边长为 \(a\) 的立方体代价为
	\(c_m a^m\)，其中 \(c_m=\omega_m(\sqrt m/2)^m\)。故
	\(\mathcal H^m_\delta(Z)\le c_m\mathcal L^m(Z)\)，特别是 Lebesgue 零测集也为 \(\mathcal H^m\) 零测集。

	给定有限测度开集 \(O\)，取所有包含于 \(O\)、直径小于 \(\delta\) 的闭球。由前面的细覆盖定理可选不交球 \(B_i\)，其并遗漏的集合为 Lebesgue 零测集，因而也为 \(\mathcal H^m\) 零测集。于是
	\[
		\mathcal H^m_\delta(O)
		\le\sum_i\omega_m r_i^m
		=\sum_i\mathcal L^m(B_i)\le\mathcal L^m(O).
	\]
	让 \(\delta\downarrow0\)，再取包含任意有限外测度集合 \(A\) 的开集，利用 Lebesgue 外正则性，得到 \(\mathcal H^m(A)\le\mathcal L^m(A)\)。外测度无穷大时先前的不等式已经给出相等。仿射平面的结论再由等距不变性得到。
\end{proof}

这里的外测度构造、归一化与五倍覆盖可对照\cite[第2--3节，第6--16页]{Simon1983}。以后同一空间中同时出现 \(\mathcal L^m\) 和 \(\mathcal H^k\)，是为了标明参数体积与不同维数的几何测度。

\subsection{测度微分与密度定理}
\label{g3:subsec:2001:03}
积分函数在一个点可以没有连续代表，但它在越来越小的球上的平均振荡仍可趋零。这种几乎处处的恢复性质将把弱导数、近似线性化和几何测度连起来。本小节的球均为 Euclid 球；开球与闭球对 Lebesgue 平均没有区别。

\begin{lemma}\label{g3:lem:measure-maximal-weak-bound}
	设 \(\sigma\) 为 \(\mathbb R^m\) 上的有限正 Borel 测度，\(m\ge1\)，定义
	\[
		M\sigma(x)=\sup_{r>0}\frac{\sigma(B_r(x))}{\omega_mr^m}.
	\]
	则对每个 \(t>0\)，
	\(\mathcal L^m\{M\sigma>t\}\le5^m\sigma(\mathbb R^m)/t\)。
	特别地，取 \(\mathrm d\sigma=|h|\,\mathrm dx\) 可得到 \(h\in L^1\) 的相同估计。
\end{lemma}
\begin{proof}
	每个满足比值大于 \(t\) 的开球都可稍微增大为仍满足该不等式的闭球，因为分母连续，测度单调。所有这些球的半径有共同上界
	\((\sigma(\mathbb R^m)/(t\omega_m))^{1/m}\)。对它们使用五倍覆盖，得到不交球 \(B_i\)，从而
	\[
		\mathcal L^m\{M\sigma>t\}
		\le5^m\sum_i\omega_mr_i^m
		<\frac{5^m}{t}\sum_i\sigma(B_i)
		\le\frac{5^m}{t}\sigma(\mathbb R^m).
	\]
	函数可测性也可直接检验：对固定 \(r\)，开球内积分随中心下半连续，故其上确界仍可测。
\end{proof}

\termidx[Lebesgue weifen dingli]{Lebesgue 微分定理}
\foreignidx{Lebesgue differentiation theorem}
\begin{theorem}[Lebesgue 微分定理]\label{g3:thm:lebesgue-differentiation}
	设 \(h\in L^1_{\mathrm{loc}}(\mathbb R^m)\)，\(m\ge1\)，选取其几乎处处有限的可测代表。则对几乎每个 \(x\)，
	\[
		\lim_{r\downarrow0}\frac1{\omega_mr^m}
		\int_{B_r(x)}|h(y)-h(x)|\,\mathrm dy=0.
	\]
	这样的点称为 \(h\) 的\term[Lebesgue dian]{Lebesgue 点}[Lebesgue point]。
\end{theorem}
\begin{proof}
	先设 \(h\in L^1(\mathbb R^m)\)。取连续紧支函数 \(v\) 逼近 \(h\)，并记左侧的上极限为 \(D h(x)\)。三角不等式与 \(v\) 的连续性给出
	\[
		D h(x)\le M(|h-v|\,\mathcal L^m)(x)+|h(x)-v(x)|.
	\]
	由前一引理与积分的 Chebyshev 不等式，
	\[
		\mathcal L^m\{D h>2t\}
		\le\frac{5^m+1}{t}\|h-v\|_{L^1}.
	\]
	连续紧支函数的 \(L^1\) 稠密性可由简单函数逼近及内外正则性得到：将有限测度集夹在紧集与开集之间，再取连接两者的连续截断函数。令逼近误差趋零，得到 \(\{D h>2t\}\) 零测；再令 \(t\) 遍历正有理数，得到 \(D h=0\) 几乎处处。局部可积情形先在每个较大球截断，较小球上的足够小平均不变，最后作可数穷尽。
\end{proof}

\begin{corollary}\label{g3:cor:lebesgue-density}
	若 \(E\subset\mathbb R^m\) 为 Lebesgue 可测集，则对几乎每个 \(x\in E\)，
	\[
		\lim_{r\downarrow0}\frac{\mathcal L^m(E\cap B_r(x))}{\omega_mr^m}=1;
	\]
	对几乎每个 \(x\notin E\)，该比值趋零。
\end{corollary}
\begin{proof}
	对 \(\mathbf1_E\in L^1_{\mathrm{loc}}\) 应用前一定理。
\end{proof}

密度点不必是内点。后面的闭集习题将给出正测度、空内部的集合，它几乎每个点仍有密度一。相应地，改变一个零测集不会改变几乎处处的密度结论。

\ProseParagraphBreak

在 \(\mathbb R^N\) 上，一个紧集测度有限、且可由紧集从内逼近的 Borel 测度称为\term[Radon cedu]{Radon 测度}[Radon measure]。\footnote{Johann Radon（1887--1956），奥地利数学家，研究测度、积分变换和凸几何。}局部有限 Borel 测度自动有这一正则性：开集可由有界闭子集递增穷尽，有限测度时内外逼近的集合族对补与可数并封闭，因此包含全部 Borel 集；一般情形再作有界穷尽。本章需要的密度记号为
\[
	\Theta^{*m}(\mu,x)=\limsup_{r\downarrow0}\frac{\mu(B_r(x))}{\omega_mr^m},
	\qquad
	\Theta_*^m(\mu,x)=\liminf_{r\downarrow0}\frac{\mu(B_r(x))}{\omega_mr^m}.
\]
上下极限相等时记共同值为 \(\Theta^m(\mu,x)\)。定义中的维数 \(m\) 可以小于环境维数 \(N\)。

\begin{proposition}\label{g3:prop:radon-lebesgue-density}
	设 \(\mu\) 为 \(\mathbb R^m\) 上的正 Radon 测度，\(m\ge1\)，其相对于 Lebesgue 测度的分解为
	\(\mu=h\mathcal L^m+\mu_s\)，其中 \(h\ge0\) 局部可积、\(\mu_s\) 集中于某个 Lebesgue 零测集。则
	\[
		\lim_{r\downarrow0}\frac{\mu(B_r(x))}{\omega_mr^m}=h(x)
		\quad\text{对 }\mathcal L^m\text{ 几乎每个 }x.
	\]
\end{proposition}
\begin{proof}
	绝对连续部分由 Lebesgue 微分定理处理。先设 \(\mu_s\) 有限，集中于零测集 \(Z\)。取紧集 \(K\subset Z\)，使 \(\mu_s(\mathbb R^m\setminus K)<\varepsilon\)，并令 \(\nu(A)=\mu_s(A\setminus K)\)。若 \(x\notin K\)，充分小的球避开 \(K\)，故
	\[
		\limsup_{r\downarrow0}\frac{\mu_s(B_r(x))}{\omega_mr^m}
		\le M\nu(x).
	\]
	上极限大于 \(t\) 的集合在 \(K\) 外的 Lebesgue 测度至多为 \(5^m\varepsilon/t\)，而 \(K\) 本身零测。令 \(\varepsilon\downarrow0\)，再遍历正有理数 \(t\)，就得到所需零极限。局部有限情形将 \(\mu_s\) 限制在逐渐增大的球内；内层球中的小尺度密度不受限制影响。
\end{proof}

这个结论以 \(\mathcal L^m\) 为“几乎处处”的基准。例如 Dirac 点质量对 Lebesgue 几乎每点的密度为零，测度本身却没有消失；在该原子处比值趋于无穷大。低维几何中的“几乎处处”也必须跟随所用的 \(\mathcal H^m\) 或加权测度。

\subsection{Hausdorff 测度与子流形体积}
\label{g3:subsec:2001:04}
现在把覆盖定义与前几章的诱导度量面积接起来。核心是一个线性事实。若 \(A:\mathbb R^m\to\mathbb R^N\) 单射，则其 \(m\) 维体积伸缩因子为
\[
	J_m(A)=\sqrt{\det(A^{\mathsf T}A)}.
\]
用正交变换把 \(A\) 写成沿 \(m\) 个方向的伸缩，Lebesgue 的线性换元与\cref{g3:thm:hausdorff-lebesgue}给出
\(\mathcal H^m(A(E))=J_m(A)\mathcal L^m(E)\)。若秩小于 \(m\)，其像在一个低维平面内；将有界片按边长 \(\delta\) 的小立方体覆盖，\(m\) 次总代价为 \(O(\delta^{m-\operatorname{rank}A})\)，故其 \(\mathcal H^m\) 测度为零。

\begin{theorem}\label{g3:thm:smooth-hausdorff-area}
	设 \(M^m\subset\mathbb R^N\) 为 \(C^1\) 嵌入子流形，\(1\le m\le N\)，带诱导度量。对每个 Borel 集 \(E\subset M\)，有
	\[
		\mathcal H^m(E)=\operatorname{Vol}_M(E).
	\]
	更具体地，若 \(F:U\subset\mathbb R^m\to M\) 是单射 \(C^1\) 参数图，且 \(\operatorname{rank}DF=m\)，则
	\begin{equation}\label{g3:eq:smooth-parametrized-area}
		\mathcal H^m(F(E))
		=\int_E\sqrt{\det(DF(x)^{\mathsf T}DF(x))}\,\mathrm dx
		\qquad(E\subset U\text{ Borel}).
	\end{equation}
\end{theorem}
\begin{proof}
	固定 \(x_0\in U\)，记 \(A=DF(x_0)\)，其最小奇异值为 \(a>0\)。由导数连续性，在以 \(x_0\) 为中心的足够小凸球上，\(\|DF(x)-A\|\le\varepsilon<a\)。沿线段积分，得到
	\[
		|F(y)-F(z)-A(y-z)|\le\varepsilon|y-z|.
	\]
	因此从 \(A\) 的像到 \(F\) 的像的对应 \(Ay\mapsto F(y)\)，及其逆对应，具有下面的距离控制：
	\[
		(1-\varepsilon/a)|A(y-z)|
		\le |F(y)-F(z)|
		\le(1+\varepsilon/a)|A(y-z)|.
	\]
	由 Lipschitz 测度不等式和线性体积公式，对该球中任意 Borel 集 \(E\)，
	\[
		(1-\varepsilon/a)^m J_m(A)\mathcal L^m(E)
		\le\mathcal H^m(F(E))
		\le(1+\varepsilon/a)^m J_m(A)\mathcal L^m(E).
	\]
	参数图是到其像的同胚，所以 \(E\mapsto\mathcal H^m(F(E))\) 为局部有限 Borel 测度；由局部 Lipschitz 控制，它对 \(\mathcal L^m\) 绝对连续。记其积分密度为 \(q\)。上式用于越来越小的 \(E=B_r(x_0)\)，再令 \(\varepsilon\downarrow0\)，说明此测度的球密度为 \(J_m(DF(x_0))\)。\Cref{g3:prop:radon-lebesgue-density}则给出 \(q=J_m(DF)\) 几乎处处，证明式\eqref{g3:eq:smooth-parametrized-area}。

	在这个参数图中，诱导度量矩阵正是 \(DF^{\mathsf T}DF\)，所以右侧就是体积密度积分。用可数个参数图覆盖 \(M\)，按先后去掉已覆盖部分，再用可数可加性，就得到整体结论。证明同时适用于带边界参数图的内部；边界局部是 \((m-1)\) 维 \(C^1\) 图，由前面的低维覆盖估计，其 \(\mathcal H^m\) 测度为零。
\end{proof}

例如 \(C^1\) 函数图 \(F(x)=(x,u(x))\) 的面积密度为
\(\sqrt{\det(I+Du^{\mathsf T}Du)}\)。标量图时化为
\(\sqrt{1+|Du|^2}\)，与\secref{g3:sec:1902}一致。这里测量的是图作为集合的大小；一个非单射参数化可能重复经过同一片区域，其参数积分须在\secref{g3:sec:2003}中加入重数后解释。

\ProseParagraphBreak

同一结论也适用于光滑黎曼环境中的光滑嵌入子流形，Hausdorff 测度采用环境距离。具体地，在一点的坐标邻域中，把该点的度量矩阵冻结，缩小邻域可使向量长度与冻结长度相差至多因子 \(1\pm\varepsilon\)。对更小邻域中的足够近点，直线段给出距离上界；离开较大坐标邻域的道路有固定正长度下界，因而不会改善这些小距离，下界则由逐段长度比较给出。故局部距离和诱导面积密度都夹在对应常系数量的 \(1\pm\varepsilon\) 次方之间。对任意小片应用上述 Euclid 结论，再令 \(\varepsilon\downarrow0\)，得到两种面积相等。

% !TeX root = ../../Book3.tex
% 纲要标题：### 20.2 Lipschitz 映射的可微性与线性化
% 纲要位置：几何学计划纲要.md，第 2933 行。
% 正文状态：初稿；尚未正式审核。

% 纲要内容开始
% 
% * 一维绝对连续函数
% * 逐方向导数和 Rademacher 定理
% * 近似微分与密度点
% * Lipschitz 映射的局部分解
% 
% 纲要内容结束

\section{Lipschitz 映射的可微性与线性化}
\label{g3:sec:2002}

Lipschitz 条件同时控制两件事：沿直线的变化可以积分，邻近方向的差商不会相差太远。前一件事产生几乎处处的方向导数，后一件事才使这些导数组成真正的线性近似。本节把这两步分开证明，再把逐点近似整理成可数个具有统一误差的可测片。

\ProseParagraphBreak
以下 \(|\cdot|\) 表示 Euclid 范数，线性映射的算子范数记为 \(\|\cdot\|\)，\(\mathcal L^m\) 表示 \(m\) 维 Lebesgue 测度。维数 \(m,n\) 均为正整数。本节的微分结果不要求二者之间有大小关系；满秩单射分片时才加上 \(m\le n\)。

\subsection{一维绝对连续函数}
\label{g3:subsec:2002:01}
\begin{definition}\label{g3:def:absolute-continuity-interval}
	设 \(a<b\)。对函数 \(u:[a,b]\to\mathbb R^n\)，若对每个 \(\varepsilon>0\)，都存在 \(\delta>0\)，使任意有限族内部两两不交的子区间 \([a_j,b_j]\subset[a,b]\) 满足
	\[
		\sum_j(b_j-a_j)<\delta
		\quad\Longrightarrow\quad
		\sum_j|u(b_j)-u(a_j)|<\varepsilon.
	\]
	则称 \(u\) \term[juedui lianxu]{绝对连续}[absolutely continuous]。这个函数类记为 \(AC([a,b];\mathbb R^n)\)。
\end{definition}

取一个区间便知绝对连续蕴含连续。若 \(u\) 为 \(L\)-Lipschitz，则左边的增量和至多为 \(L\sum_j(b_j-a_j)\)，所以 Lipschitz 蕴含绝对连续。下面的定理说明，这种对许多短区间的共同控制恰好排除了不能由积分恢复的变化。

\begin{theorem}[绝对连续函数的微积分基本定理]\label{g3:thm:ac-fundamental-calculus}
	设 \(a<b\)，\(u:[a,b]\to\mathbb R^n\)。下列两件事等价：
	\begin{enumerate}
		\item \(u\) 绝对连续；
		\item 存在 \(g\in L^1((a,b);\mathbb R^n)\)，使每个 \(t\in[a,b]\) 都满足
		\begin{equation}\label{g3:eq:ac-integral-representation}
			u(t)=u(a)+\int_a^t g(s)\,\mathrm ds.
		\end{equation}
	\end{enumerate}
	在此情形，\(u\) 几乎处处可微，且 \(u'=g\) 几乎处处。特别地，\(u\) 为 \(L\)-Lipschitz 当且仅当它绝对连续且 \(|u'|\le L\) 几乎处处。
\end{theorem}
\begin{proof}
	先证实值情形。设 \(u\) 绝对连续。选取对应于 \(\varepsilon=1\) 的 \(\delta\)，把 \([a,b]\) 分成有限个长度小于 \(\delta\) 的区间。任一分割再插入这些固定分点后，在每一个小区间内的增量绝对值之和至多为一。因此总变差
	\[
		V(t)=\sup_{a=t_0<\cdots<t_k=t}
		\sum_{i=1}^k|u(t_i)-u(t_{i-1})|
	\]
	有限。细分分割并取上确界可得
	\(V(t)-V(s)=\operatorname{Var}(u;[s,t])\)。对任意有限族内部不交的小区间，分别用有限分割逼近其中的变差；这些分割的区间总长度仍等于原来的总长度。绝对连续性于是给出
	\[
		\sum_j\bigl(V(b_j)-V(a_j)\bigr)<\varepsilon
	\]
	（必要时先对 \(\varepsilon/2\) 选择 \(\delta\)，再取上确界）。故 \(V\) 也绝对连续。

	函数
	\[
		v_\pm(t)=\frac{V(t)\pm(u(t)-u(a))}{2}
	\]
	均单调不减且绝对连续，并且 \(u-u(a)=v_+-v_-\)。对连续单调函数 \(v\)，区间增量
	\(\nu((s,t])=v(t)-v(s)\) 定义一个有限正 Borel 测度。这里可直接使用区间测度的延拓构造：增量的有限可加性来自相消求和；可数区间覆盖的次可加性，先把左端略向右移、把覆盖区间略放大，再对所得紧区间取有限子覆盖，最后用 \(v\) 的连续性取极限。因而该区间函数是预测度，其外测度延拓具有所给区间值；另置 \(\nu(\{a\})=0\)。

	当 \(v\) 绝对连续时，这个测度满足 \(\nu\ll\mathcal L^1\)。事实上，把任意 Lebesgue 零测集盖在总长度小于 \(\delta\) 的开集中，将开集写成两两不交的开区间；对有限部分使用绝对连续性，再取单调极限，即得其 \(\nu\)-测度至多为任意给定的 \(\varepsilon\)。端点没有原子，因为 \(v\) 连续。应用本章沿用的绝对连续有限测度的积分密度表示，得到 \(g_\pm\in L^1\)，使
	\[
		v_\pm(t)-v_\pm(a)=\int_a^t g_\pm(s)\,\mathrm ds.
	\]
	相减便得 \eqref{g3:eq:ac-integral-representation}。

	反之，若有该积分表示，则
	\[
		\sum_j|u(b_j)-u(a_j)|
		\le\int_{\bigcup_j(a_j,b_j)}|g(s)|\,\mathrm ds.
	\]
	\(L^1\) 积分对小测度集的绝对连续性给出所需结论。由\cref{g3:thm:lebesgue-differentiation}，几乎处处的 \(t\) 满足
	\[
		\frac1{2r}\int_{t-r}^{t+r}|g(s)-g(t)|\,\mathrm ds\longrightarrow0.
	\]
	单侧区间上的平均误差至多为相应对称区间平均误差的两倍，故
	\[
		\frac{u(t+h)-u(t)}h-g(t)
		=\frac1h\int_t^{t+h}(g(s)-g(t))\,\mathrm ds\longrightarrow0.
	\]
	这证明了几乎处处的导数公式。

	向量值情形逐分量应用刚才的证明；有限个分量的绝对连续性与向量形式等价。若 \(u\) 为 \(L\)-Lipschitz，其存在的导数必有模长至多 \(L\)。反之，由积分表示和 \(|g|\le L\)，立即得到 \(|u(t)-u(s)|\le L|t-s|\)。
\end{proof}

\begin{corollary}\label{g3:cor:ac-integration-parts}
	设 \(u\in AC([a,b];\mathbb R^n)\)，\(\varphi\in C^1([a,b];\mathbb R)\)。则
	\[
		\int_a^b u'(t)\varphi(t)\,\mathrm dt
		=[u\varphi]_a^b-\int_a^b u(t)\varphi'(t)\,\mathrm dt.
	\]
\end{corollary}
\begin{proof}
	\(u\) 与 \(\varphi\) 在紧区间上有界，乘积增量的估计说明 \(u\varphi\) 绝对连续。在 \(u\) 可微的满测集上使用乘积求导法则，再对 \(u\varphi\) 应用\cref{g3:thm:ac-fundamental-calculus}即可。
\end{proof}

\begin{example}\label{g3:exm:ac-lipschitz-distinction}
	\(|t|\) 在 \([-1,1]\) 上 Lipschitz，却在原点不可微。\(\sqrt t\) 在 \([0,1]\) 上绝对连续，因为它是可积函数 \(1/(2\sqrt t)\) 的原函数，却不是 Lipschitz。这两个例子分别说明：Lipschitz 条件允许零测集上的尖点，绝对连续性则还允许无界但可积的导数。
\end{example}

\subsection{逐方向导数和 Rademacher 定理}
\label{g3:subsec:2002:02}
对一个固定方向，上一小节已经给出几乎处处的导数。困难在于：每一个方向原先都带着自己的例外集，而即使所有方向导数同时存在，它们也未必组成线性映射。Rademacher 定理\footnote{Hans Rademacher（1892--1969），德国出生的数学家，后在美国任教，研究实变函数与解析数论，建立了 Lipschitz 映射的几乎处处可微性定理。}的下述证明先用分布的线性性解决后一问题，再只取可数个方向来控制例外集。这个证明可与 \cite[第5节，定理5.2，第30--32页]{Simon1983} 对照。

\termidx[Rademacher dingli]{Rademacher 定理}
\foreignidx{Rademacher theorem}
\begin{theorem}[Rademacher 定理]\label{g3:thm:rademacher}
	设 \(m,n\ge1\)，\(f:\mathbb R^m\to\mathbb R^n\) 为 \(L\)-Lipschitz 映射。则几乎每个 \(x\) 都存在一个线性映射 \(\mathrm Df(x):\mathbb R^m\to\mathbb R^n\)，使
	\begin{equation}\label{g3:eq:rademacher-total-differential}
		\lim_{h\to0}\frac{|f(x+h)-f(x)-\mathrm Df(x)h|}{|h|}=0.
	\end{equation}
	并且 \(\|\mathrm Df(x)\|\le L\)。在可微的满测集上，\(\mathrm Df\) 的各列就是 \(f\) 的弱偏导数。
\end{theorem}
\begin{proof}
	固定单位向量 \(v\)。每条平行于 \(v\) 的直线上，函数 \(t\mapsto f(z+tv)\) 都是 \(L\)-Lipschitz，因此在每个有界区间上绝对连续。由\cref{g3:thm:ac-fundamental-calculus} 和 Fubini 定理，方向导数
	\[
		g_v(x)=\lim_{t\to0}\frac{f(x+tv)-f(x)}t
	\]
	对几乎每个 \(x\) 存在，且 \(|g_v|\le L\)。其存在集可测：因差商关于非零 \(t\) 连续，极限的 Cauchy 条件只须在有理 \(t\) 上检验。把 \(g_v\) 在例外集上置为零，得到有界可测函数。

	沿平行直线应用\cref{g3:cor:ac-integration-parts}，再用 Fubini 定理，可得每个 \(\varphi\in C_c^\infty(\mathbb R^m)\) 满足
	\begin{equation}\label{g3:eq:directional-weak-linearity}
		\int_{\mathbb R^m}g_v\varphi\,\mathrm dx
		=-\int_{\mathbb R^m}f\,\partial_v\varphi\,\mathrm dx
		=\sum_{i=1}^m v_i\int_{\mathbb R^m}g_{e_i}\varphi\,\mathrm dx.
	\end{equation}
	因此
	\(g_v=\sum_i v_i g_{e_i}\) 几乎处处。这里使用的是局部可积函数所定义分布的唯一性：若 \(w\in L^1_{\mathrm{loc}}\) 对所有测试函数的积分为零，则在任意内缩开集上 \(w\) 与光滑磨光核的卷积为零；卷积在 \(L^1_{\mathrm{loc}}\) 中趋向 \(w\)，故 \(w=0\) 几乎处处。这与第十一章弱导数的唯一性论证相同。

	取单位球面上的可数稠密集 \(\mathcal V\)，使它包含 \(e_1,\ldots,e_m\)。去掉各 \(v\in\mathcal V\) 的导数不存在集和 \eqref{g3:eq:directional-weak-linearity} 所得等式的例外集；其并仍为零测集。对剩下的每个 \(x\)，定义
	\[
		A_x h=\sum_{i=1}^m h_i g_{e_i}(x).
	\]
	于是所有 \(v\in\mathcal V\) 同时满足 \(g_v(x)=A_xv\)。又因 \(|g_v(x)|\le L\)，稠密性给出 \(\|A_x\|\le L\)。

	固定这样的 \(x\) 和 \(\varepsilon>0\)。若 \(L=0\)，结论显然；以下设 \(L>0\)。从 \(\mathcal V\) 取有限个向量 \(v_1,\ldots,v_q\)，构成单位球面的 \(\varepsilon/(4L)\)-网。有限个方向导数的存在性给出 \(r_0>0\)，使
	\[
		|f(x+rv_j)-f(x)-rA_xv_j|\le\frac\varepsilon2r
		\qquad(0<r<r_0,\ 1\le j\le q).
	\]
	若 \(h=r\theta\)，\(|\theta|=1\)，选取 \(|\theta-v_j|<\varepsilon/(4L)\)，则
	\begin{align*}
		|f(x+h)-f(x)-A_xh|
		&\le |f(x+r\theta)-f(x+rv_j)|
		 +|f(x+rv_j)-f(x)-rA_xv_j|\\
		&\qquad+r|A_x(v_j-\theta)|\\
		&\le 2Lr|\theta-v_j|+\frac\varepsilon2r
		\le\varepsilon r.
	\end{align*}
	这就是全微分，而不仅是逐方向极限。最后，取 \(v=e_i\) 的积分恒等式说明这些导数是弱偏导数。
\end{proof}

Rademacher 定理也适用于开集上的局部 Lipschitz 映射：用可数个 Lipschitz 常数有限的球覆盖开集，在每个球内作同样的直线分部积分和有限方向网论证即可。所用测试函数紧含于球，差商也只取足够小的增量。

\ProseParagraphBreak
证明中的顺序值得保留：先得到弱导数，再让可数个方向导数与同一个矩阵相合，最后才得到任意方向上的统一余项。如果删去中间的线性关系，“每个方向都可微”本身不足以推出全微分。

\subsection{近似微分与密度点}
\label{g3:subsec:2002:03}
光滑几何常在开集上求微分；可测几何却需要在只占据局部几乎全部体积的集合上识别线性近似。密度点提供了这样的比较位置。

\begin{definition}\label{g3:def:approximate-differential}
	设 \(E\subset\mathbb R^m\) 为 Lebesgue 可测集，\(f:E\to\mathbb R^n\) 可测，\(x\in E\) 是 \(E\) 的密度为一的点。若线性映射 \(A:\mathbb R^m\to\mathbb R^n\) 对每个 \(\varepsilon>0\) 都满足
	\begin{equation}\label{g3:eq:approximate-differential}
		\lim_{r\downarrow0}
		\frac{\mathcal L^m\{y\in E\cap B(x,r):
		 |f(y)-f(x)-A(y-x)|>\varepsilon|y-x|\}}
		 {\omega_m r^m}=0,
	\end{equation}
	就称 \(f\) 在 \(x\) \term[jinsi weifen]{近似可微}[approximately differentiable]，并把 \(A\) 称为其近似微分，记为 \(\operatorname{ap}\mathrm Df(x)\)。
\end{definition}
\symidx{\(\operatorname{ap}\mathrm Df(x)\)：可测映射在密度点的近似微分}

\begin{lemma}\label{g3:lem:density-no-large-holes}
	设 \(E\subset\mathbb R^m\) 可测，\(x\) 是其密度为一的点。则
	\begin{equation}\label{g3:eq:density-no-large-holes}
		\sup_{|z-x|\le r}\operatorname{dist}(z,E)=o(r)
		\qquad(r\downarrow0).
	\end{equation}
	若可测映射 \(f:E\to\mathbb R^n\) 在 \(x\in E\) 存在近似微分，则它唯一。
\end{lemma}
\begin{proof}
	若第一式不成立，则存在 \(\eta>0\)、\(r_j\downarrow0\) 和 \(z_j\in\overline B(x,r_j)\)，使 \(\operatorname{dist}(z_j,E)>\eta r_j\)。可设 \(\eta<1\)。于是 \(B(z_j,\eta r_j)\subset B(x,2r_j)\setminus E\)，后者占 \(B(x,2r_j)\) 的比例至少为 \((\eta/2)^m\)，与密度为一矛盾。

	设 \(A,B\) 都是近似微分，而 \(A-B\ne0\)。存在单位向量 \(v\)、\(c>0\) 和 \(v\) 的一个球面邻域，使其中的 \(\theta\) 均满足 \(|(A-B)\theta|>c\)。该方向锥在每个 \(B(x,r)\) 内占一个固定的正体积比例。取 \(\varepsilon<c/3\)。在 \(E\) 中除去两个近似展开的坏集后，三角不等式却给出 \(|(A-B)(y-x)|\le2\varepsilon|y-x|\)。两个坏集与 \(E\) 的补集在 \(x\) 的密度都为零，无法覆盖上述方向锥，矛盾。
\end{proof}

\begin{proposition}\label{g3:prop:lipschitz-extension-density}
	设 \(E\subset\mathbb R^m\) 非空，\(f:E\to\mathbb R^n\) 为 \(L\)-Lipschitz 映射。则存在全空间上的 Lipschitz 延拓 \(F:\mathbb R^m\to\mathbb R^n\)，其 Lipschitz 常数至多为 \(\sqrt nL\)。若 \(E\) 可测，则有以下结论。
	\begin{enumerate}
		\item 在 \(E\) 的密度为一的点 \(x\in E\)，只要一个 Lipschitz 延拓可微，所有 Lipschitz 延拓都可微，且导数相同。
		\item 若 \(f\) 在这样的点近似可微，则任一 Lipschitz 延拓都在该点可微，导数等于 \(\operatorname{ap}\mathrm Df(x)\)。反之，延拓的可微性蕴含 \(f\) 的近似可微性。
		\item 因而在 \(E\) 的几乎每一点，存在不依赖延拓的微分，且它等于近似微分，其算子范数至多为原映射的 Lipschitz 常数 \(L\)。特别地，两个 Lipschitz 映射若在 \(E\) 上相等，它们的微分在 \(E\) 上几乎处处相等。
	\end{enumerate}
\end{proposition}
\begin{proof}
	对每个分量定义
	\[
		F_i(z)=\inf_{a\in E}\{f_i(a)+L|z-a|\}.
	\]
	固定 \(a_0\in E\)，由 Lipschitz 条件和三角不等式，下确界中的每一项均不小于 \(f_i(a_0)-L|z-a_0|\)，而取 \(a=a_0\) 给出有限上界；所以此下确界有限。对 \(z\in E\) 它等于 \(f_i(z)\)，且三角不等式给出 \(|F_i(z)-F_i(w)|\le L|z-w|\)。合并分量便得所述延拓。这一标量延拓构造也见 \cite[第5节，定理5.1，第29--30页]{Simon1983}。

	若 \(F,G\) 是两个 Lipschitz 延拓，常数分别为 \(L_F,L_G\)，则
	\[
		|F(z)-G(z)|\le(L_F+L_G)\operatorname{dist}(z,E).
	\]
	由\cref{g3:lem:density-no-large-holes}，在密度为一的 \(x\in E\) 处，右边为 \(o(|z-x|)\)。而 \(F(x)=G(x)\)，所以其中一个的一阶展开也适用于另一个。

	若 \(A=\operatorname{ap}\mathrm Df(x)\)，固定 \(\varepsilon>0\)，令
	\[
		E_\varepsilon=\{y\in E:
		 |f(y)-f(x)-A(y-x)|\le\varepsilon|y-x|\}.
	\]
	这是在 \(x\) 密度为一的集合。对 \(z\to x\)，由无大空洞估计可选 \(y\in E_\varepsilon\)，使 \(|z-y|=o(|z-x|)\)。于是
	\[
		|F(z)-F(x)-A(z-x)|
		\le(L_F+\|A\|)|z-y|+\varepsilon|y-x|.
	\]
	除以 \(|z-x|\) 后取上极限，所得上界至多为 \(\varepsilon\)。再令 \(\varepsilon\downarrow0\)，便得 \(\mathrm DF(x)=A\)。反向蕴含由定义直接得到。

	最后，\cref{g3:cor:lebesgue-density,g3:thm:rademacher} 分别排除非密度点和延拓的不可微点；这些集合都为零测集。在其余的点 \(x\)，对任意 \(v\) 用无大空洞估计选取 \(y_t\in E\)，使 \(y_t=x+tv+o(|t|)\)。把 \(|f(y_t)-f(x)|\le L|y_t-x|\) 除以 \(|t|\) 并令 \(t\to0\)，得到 \(|\mathrm DF(x)v|\le L|v|\)；所以导数界保留了原片上的常数。对两个相等的映射，使用同一限制的延拓独立性即可。
\end{proof}

\begin{example}\label{g3:exm:approximate-not-classical}
	函数 \(f=\mathbf1_{\mathbb Q}:\mathbb R\to\mathbb R\) 在每个无理点的近似微分为零：取 \(A=0\)，坏集包含在零测集 \(\mathbb Q\) 中。但它在每一点都不连续，因而处处没有通常微分。在有理点，所给点值为一，反而不存在近似微分。这也说明改变零测集上的点值不影响“几乎处处近似可微”，却可能改变被修改点本身的结论。
\end{example}

由上述命题，一个 Lipschitz 映射在可测集上恒定时，其微分在该集合上几乎处处为零。这里集合可以没有内点；起作用的是密度为一，而不是开邻域。

\subsection{Lipschitz 映射的局部分解}
\label{g3:subsec:2002:04}
下一步不再扩大可微点集，而是把不同点各自拥有的“小尺度”整理为共同尺度。下面的分片同时冻结导数矩阵和 Taylor 余项；面积与余面积公式正是在这些片上与线性映射比较。

\begin{theorem}\label{g3:thm:lipschitz-linearization-pieces}
	设 \(f:\mathbb R^m\to\mathbb R^n\) 为 Lipschitz 映射，\(D\) 为其全体可微点。则 \(D\) 是 Borel 集，\(x\mapsto\mathrm Df(x)\) 在 \(D\) 上 Borel 可测。对每个 \(\varepsilon>0\)，存在两两不交的可数族 Borel 集 \(E_j\)，以及线性映射 \(A_j:\mathbb R^m\to\mathbb R^n\)，使 \(D=\bigcup_jE_j\)，且
	\begin{align}
		\|\mathrm Df(x)-A_j\|&\le\varepsilon
		&& (x\in E_j),\label{g3:eq:piece-derivative-close}\\
		|f(y)-f(x)-A_j(y-x)|&\le\varepsilon|y-x|
		&& (x,y\in E_j).\label{g3:eq:piece-uniform-linearization}
	\end{align}
\end{theorem}
\begin{proof}
	首先把可测性具体写出。令 \(D_0\) 为每列差商
	\(k(f(x+e_i/k)-f(x))\) 都收敛的点集；Cauchy 条件说明 \(D_0\) 是 Borel 集。在其上把这些极限组成矩阵 \(A(x)\)，在其外置 \(A(x)=0\)，得到 Borel 矩阵函数。对正整数 \(p,q\)，条件
	\[
		|f(x+h)-f(x)-A(x)h|\le p^{-1}|h|
		\quad\text{对所有 }h\in\mathbb Q^m,\quad0<|h|<q^{-1}
	\]
	定义一个 Borel 点集 \(D_{pq}\)。关于 \(h\) 的连续性使同一不等式对所有实增量成立。于是
	\[
		D=D_0\cap\bigcap_{p=1}^\infty\bigcup_{q=1}^\infty D_{pq},
		\qquad \mathrm Df=A\ \text{于 }D,
	\]
	证明了第一部分。

	取矩阵空间中的可数稠密族 \(\{A_j\}\)。给定 \(\varepsilon>0\)，把 \(D\) 按以下两个条件分组：
	\[
		\|\mathrm Df(x)-A_j\|\le\varepsilon/2,
	\]
	以及
	\[
		|f(x+h)-f(x)-\mathrm Df(x)h|
		\le(\varepsilon/2)|h|
		\quad(0<|h|<q^{-1}).
	\]
	后一个条件同样只须在有理增量上检验，故各组均为 Borel 集；可微性说明这些组覆盖 \(D\)。再将每组与直径严格小于 \(q^{-1}\) 的可数网格小方块相交。若 \(x,y\) 在同一交片中，把 Taylor 估计用于以 \(x\) 为基点的增量 \(y-x\)，便得
	\[
		|f(y)-f(x)-A_j(y-x)|
		\le\frac\varepsilon2|y-x|
		 +\|\mathrm Df(x)-A_j\|\,|y-x|
		\le\varepsilon|y-x|.
	\]
	最后把所得可数覆盖列成一列，并从每片扣除此前各片，就得到两两不交的 Borel 分解；两个估计在子集上仍成立。
\end{proof}

\begin{corollary}\label{g3:cor:full-rank-bilipschitz-pieces}
	设 \(1\le m\le n\)，\(f:\mathbb R^m\to\mathbb R^n\) 为 \(L\)-Lipschitz 映射。其可微且 \(\operatorname{rank}\mathrm Df=m\) 的点集可分成可数个两两不交的 Borel 片。在每一片 \(E\) 上，存在 \(c_E>0\)，使
	\[
		c_E|x-y|\le|f(x)-f(y)|\le L|x-y|
		\qquad(x,y\in E).
	\]
	特别地，\(f|_E\) 单射，其到像集的逆映射也是 Lipschitz。
\end{corollary}
\begin{proof}
	令 \(\sigma_{\min}(A)=\min_{|v|=1}|Av|\)。它对矩阵连续，且 \(\operatorname{rank}A=m\) 当且仅当 \(\sigma_{\min}(A)>0\)。把满秩集写成 Borel 集
	\(D_k=\{x\in D:\sigma_{\min}(\mathrm Df(x))\ge1/k\}\) 的并。对每个 \(k\)，以 \(\varepsilon=1/(4k)\) 应用\cref{g3:thm:lipschitz-linearization-pieces}，再与 \(D_k\) 相交。在这样的片上，
	\[
		\sigma_{\min}(A_j)\ge\frac1k-\varepsilon,
		\qquad
		|f(y)-f(x)|\ge\left(\frac1k-2\varepsilon\right)|y-x|
		=\frac1{2k}|y-x|.
	\]
	上界来自原映射的 Lipschitz 条件。对全部片再作逐次去重即可。
\end{proof}

\begin{corollary}\label{g3:cor:compact-linearization-pieces}
	设 \(f:\mathbb R^m\to\mathbb R^n\) 为 Lipschitz 映射，\(m,n\ge1\)。取其可微点集的一致线性化分解，或在 \(m\le n\) 时取满秩可微点集的双 Lipschitz 分解。
	在\cref{g3:thm:lipschitz-linearization-pieces} 或\cref{g3:cor:full-rank-bilipschitz-pieces} 的任一分解中，各片可用可数个紧子集覆盖，至多剩下一个 \(\mathcal L^m\)-零测集。因而可以得到两两不交的 \(\sigma\)-紧片 \(P_j\)，覆盖原来的点集至零测余集，每片仍满足原来的估计，并且 \(f(P_j)\) 是 Borel 集。
\end{corollary}
\begin{proof}
	对每片先与有界球相交，再由 Lebesgue 测度的内正则性选出紧子集，使被遗漏测度依次趋于零。对所有片和球取并，得到紧集列 \(K_1,K_2,\ldots\)，未覆盖部分为零测集。令
	\[
		P_j=K_j\setminus\bigcup_{i<j}K_i.
	\]
	有限个紧集的并是闭集，故 \(P_j\) 是紧集 \(K_j\) 的相对开子集；用到该闭集距离至少为 \(1/r\) 的紧集穷尽它，可知它是 \(\sigma\)-紧的。\(j=1\) 时直接取 \(P_1=K_1\)。每个 \(P_j\) 含于其原来的片，且不同 \(P_j\) 不交。连续映射把紧集变成紧集，因此 \(f(P_j)\) 也是可数个紧集的并，特别是 Borel 集。
\end{proof}

当映射起初只定义在一个可测集上时，先使用\cref{g3:prop:lipschitz-extension-density} 延拓，再把这些片与原集合相交。若需要紧片，就直接对交集使用内正则性。由密度点处的延拓独立性，所得到的几乎处处微分是原映射自身的数据。

\ProseParagraphBreak
满秩假设在这里承担明确的作用：线性化矩阵的最小奇异值给出反向距离估计。若 \(m\ge n\) 且 \(\mathrm Df\) 的秩为 \(n\)，则选择适当的 \(m-n\) 个源坐标 \(z\)，可使 \(G=(f,z):\mathbb R^m\to\mathbb R^m\) 的微分可逆；有限种坐标选择和上述分片就把这部分变成等维双 Lipschitz 坐标。这将是余面积公式中沿纤维积分的出发点。
% 纲要条目排为编号小节标题，并进入目录与书签；说明段仍保留于纲要注释中。

% !TeX root = ../../Book3.tex
% 纲要标题：### 20.3 面积公式
% 纲要位置：几何学计划纲要.md，第 2940 行。
% 正文状态：初稿；尚未正式审核。

% 纲要内容开始
% 
% * 单射 Lipschitz 片的 Jacobian
% * 重数函数和可测性
% * 由线性模型到一般映射
% * 非单射情形的分片求和
% * 几何面积与参数面积的区别
% 
% 纲要内容结束

\section{面积公式}
\label{g3:sec:2003}

一个参数片可以把同一段曲面描画多次。参数域上的面积积分会保留这些重复，而像集的 Hausdorff 测度只计算每个像点一次。面积公式把二者联系起来：微分给出每个参数点的面积放大率，重数则记录同一像点有多少个原像。光滑情形及其带重数、带权形式可参看\cite[第8节，式(8.1)--(8.5)]{Simon1983}。本节从线性映射和上一节的 Lipschitz 分片出发，证明任意可测参数集上的公式。

\subsection{单射 Lipschitz 片的 Jacobian}
\label{g3:subsec:2003:01}
以下 \(1\le m\le n\)，\(\mathcal L^m\) 表示 \(\mathbb R^m\) 上的 Lebesgue 测度，\(\mathcal H^m\) 采用\subsecref{g3:subsec:2001:01}的 \(\omega_m/2^m\) 归一化。测度记号用于尚未证明可测的集合时，先指相应外测度。\Cref{g3:thm:hausdorff-lebesgue}保证一个 \(m\) 维 Euclid 平面上的 \(\mathcal H^m\) 就是它的通常体积。

\begin{definition}\label{g3:def:area-jacobian}
	对线性映射 \(T:\mathbb R^m\to\mathbb R^n\)，定义其 \(m\) 维\term[Jacobian mianjiyinzi]{Jacobian 面积因子}[area Jacobian]为
	\[
		J_m(T)\coloneqq\sqrt{\det(T^{\mathsf T}T)}.
	\]
	对 Lipschitz 映射 \(f:\mathbb R^m\to\mathbb R^n\)，在可微点定义 \(J_m f(x)=J_m(\mathrm Df(x))\)，在不可微点令其为零。不可微集为 \(\mathcal L^m\) 零测集，故这一补充值不影响积分。
\end{definition}

矩阵 \(T^{\mathsf T}T\) 的元素是 \(\langle Te_i,Te_j\rangle\)，所以 \(J_m(T)\) 是这 \(m\) 个向量所张平行多面体的体积。若奇异值为 \(s_1,\ldots,s_m\ge0\)，则 \(J_m(T)=s_1\cdots s_m\)；它为零当且仅当 \(\operatorname{rank}T<m\)。特别地，\(J_1(T)=|Te_1|\)，而 \(m=n\) 时 \(J_m(T)=|\det T|\)。上一节证明微分矩阵在可微集上可测，故 \(J_mf\) 可测，且 \(0\le J_mf\le(\operatorname{Lip}f)^m\)。

\begin{lemma}\label{g3:lem:linear-area-formula}
	设 \(1\le m\le n\)，\(T:\mathbb R^m\to\mathbb R^n\) 为线性映射。对任意集合 \(A\subset\mathbb R^m\)，有外测度恒等式
	\begin{equation}\label{g3:eq:linear-area-formula}
		\mathcal H^m(T(A))=J_m(T)\mathcal L^m(A),
	\end{equation}
	其中 \(0\cdot\infty\) 按零解释。
\end{lemma}
\begin{proof}
	先设 \(T\) 满秩。对称矩阵 \(T^{\mathsf T}T\) 可由正交基对角化；相应的 \(Te_i/s_i\) 是像平面的正交归一基。因此在源、靶各作一次等距变换后，\(T\) 化为
	\[
		(x_1,\ldots,x_m)\longmapsto(s_1x_1,\ldots,s_mx_m,0,\ldots,0),\qquad s_i>0.
	\]
	Hausdorff 外测度在等距变换下不变。坐标伸缩使轴平行长方体的体积乘以 \(s_1\cdots s_m\)；用可数长方体覆盖计算 Lebesgue 外测度，得到伸缩后的外测度至多为这一倍。再对逆伸缩使用相同不等式，便得到等号。结合平面上 \(\mathcal H^m=\mathcal L^m\)，即得所述公式。

	若 \(\operatorname{rank}T=k<m\)，其像包含在一个 \(k\) 维平面中。该平面的任一有界部分可由至多 \(C\delta^{-k}\) 个直径不超过 \(C\delta\) 的小立方体覆盖，覆盖的 \(m\) 维代价不超过 \(C\delta^{m-k}\)，随 \(\delta\downarrow0\) 趋零。对有界部分作可数并，得整个像平面的 \(\mathcal H^m\) 为零；同时 \(J_m(T)=0\)。
\end{proof}

线性公式说明面积因子应取什么值。下面的比较说明，若在一个片上控制了两个点之间的误差，那么面积也随之受到控制。这比只固定一个点写 Taylor 展开更适于分片求和。

\begin{lemma}\label{g3:lem:near-linear-area-comparison}
	设 \(1\le m\le n\)，\(T:\mathbb R^m\to\mathbb R^n\) 为满秩线性映射，且 \(|Tv|\ge\alpha|v|\)、\(\alpha>0\)。设 \(E\subset\mathbb R^m\)，\(f:E\to\mathbb R^n\) 满足
	\[
		|f(y)-f(x)-T(y-x)|\le\varepsilon|y-x|\quad(x,y\in E),\qquad 0\le\varepsilon<\alpha.
	\]
	记 \(\rho=\varepsilon/\alpha\)。则 \(f\) 单射，而且对任意 \(A\subset E\)，
	\begin{equation}\label{g3:eq:near-linear-area-comparison}
		(1-\rho)^mJ_m(T)\mathcal L^m(A)
		\le\mathcal H^m(f(A))
		\le(1+\rho)^mJ_m(T)\mathcal L^m(A).
	\end{equation}
	此外，若线性映射 \(B:\mathbb R^m\to\mathbb R^n\) 满足 \(\lVert B-T\rVert\le\varepsilon\)，则
	\begin{equation}\label{g3:eq:near-linear-jacobian-comparison}
		(1-\rho)^mJ_m(T)\le J_m(B)\le(1+\rho)^mJ_m(T).
	\end{equation}
\end{lemma}
\begin{proof}
	把 \(f\) 看作定义在 \(T(E)\) 上的映射 \(h=f\circ(T|_E)^{-1}\)。由假设，
	\[
		(1-\rho)|Tx-Ty|\le|f(x)-f(y)|\le(1+\rho)|Tx-Ty|.
	\]
	所以 \(h\) 的 Lipschitz 常数不超过 \(1+\rho\)，其逆的常数不超过 \((1-\rho)^{-1}\)。对这两个映射分别使用\cref{g3:prop:hausdorff-lipschitz}，再代入\cref{g3:lem:linear-area-formula}，即得式\eqref{g3:eq:near-linear-area-comparison}。

	对线性映射 \(B\)，同样有 \((1-\rho)|Tv|\le|Bv|\le(1+\rho)|Tv|\)。令 \(S=(T^{\mathsf T}T)^{1/2}\)，则 \(|Tv|=|Sv|\)。在 \(v=S^{-1}w\) 上应用不等式，说明对称矩阵 \(S^{-1}B^{\mathsf T}BS^{-1}\) 的全部特征值位于 \([(1-\rho)^2,(1+\rho)^2]\)。取行列式再开平方，并用 \(\det S=J_m(T)\)，得到式\eqref{g3:eq:near-linear-jacobian-comparison}。
\end{proof}

\subsection{重数函数和可测性}
\label{g3:subsec:2003:02}
设 \(f:\mathbb R^m\to\mathbb R^n\) 为 Lipschitz 映射，\(A\subset\mathbb R^m\) 可测。定义像点 \(y\) 的\term[chongshu hanshu]{重数函数}[multiplicity function]为
\[
	N(f,A,y)\coloneqq\mathcal H^0\bigl(A\cap f^{-1}(y)\bigr).
\]
这里 \(\mathcal H^0\) 是计数测度：有限纤维取其点数，任意无限纤维取 \(+\infty\)。即使 \(A\) 紧，重数也未必有限，例如常值映射把整个参数集压到一点。

\begin{proposition}\label{g3:prop:area-multiplicity-measurable}
	设 \(1\le m\le n\)，\(f:\mathbb R^m\to\mathbb R^n\) 为 Lipschitz 映射，\(A\subset\mathbb R^m\) 为 Lebesgue 可测集。则 \(f(A)\) 为 \(\mathcal H^m\) 可测集，\(N(f,A,\cdot)\) 为非负 \(\mathcal H^m\) 可测函数。更一般地，对任意非负可测函数 \(a:A\to[0,\infty]\)，函数
	\[
		S_a(y)\coloneqq\sum_{x\in A\cap f^{-1}(y)}a(x)
	\]
	也是 \(\mathcal H^m\) 可测的；任意指标集上的非负和均解释为所有有限子和的上确界。
\end{proposition}
\begin{proof}
	先设 \(A=K\) 紧。对正整数 \(q,l\)，考虑 \(K^q\) 的闭子集
	\[
		C_{q,l}=\{(x_1,\ldots,x_q):f(x_1)=\cdots=f(x_q),\ |x_i-x_j|\ge l^{-1}\ (i\ne j)\}.
	\]
	它是紧集，故其在连续映射 \((x_1,\ldots,x_q)\mapsto f(x_1)\) 下的像 \(Y_{q,l}\) 紧。有限个不同的原像点之间有正的最小距离，因此
	\[
		\{y:N(f,K,y)\ge q\}=\bigcup_{l=1}^{\infty}Y_{q,l}.
	\]
	这证明紧集上的重数为 Borel 函数。

	由 Lebesgue 测度的内正则性，可取递增紧集 \(K_j\subset A\)，使 \(Z=A\setminus\bigcup_jK_j\) 为零测集。对于无界 \(A\)，先在 \(A\cap[-r,r]^m\) 中用紧集作任意精度的内逼近，再将得到的可数个紧集合并为递增列，即可得到同样结论。由\cref{g3:prop:hausdorff-lipschitz,g3:thm:hausdorff-lebesgue}，\(\mathcal H^m(f(Z))=0\)。而
	\[
		f(A)=\bigcup_jf(K_j)\ \cup\ f(Z),
	\]
	右端是一个 Borel 集与零测集的并，因而可测。对 \(y\notin f(Z)\)，任意有限组原像点均包含于某个 \(K_j\)，所以
	\[
		N(f,A,y)=\lim_{j\to\infty}N(f,K_j,y).
	\]
	Hausdorff 测度完备，故在零测集 \(f(Z)\) 上取实际重数仍得到可测函数。

	若 \(a\) 是非负简单函数，把它写成有限个互不交可测集的示性函数的线性组合，\(S_a\) 就是相应重数函数的同一线性组合。一般情形取非负简单函数 \(a_j\uparrow a\)。有限子和与递增极限可交换：一方面 \(S_{a_j}\le S_a\)，另一方面任意有限原像集合上的和都是这些递增和的极限。取所有有限集合的上确界，得 \(S_{a_j}\uparrow S_a\)。于是 \(S_a\) 可测。
\end{proof}

这项可测性论证只用了紧集、零测集的像和非负极限。尤其在 \(m<n\) 时，无须假设整个 \(\mathbb R^n\) 上的 \(\mathcal H^m\) 为 \(\sigma\) 有限测度。后面在目标空间所作的是可测函数列的单调收敛。

\subsection{由线性模型到一般映射}
\label{g3:subsec:2003:03}
微分满秩时，上一节已把映射分成可数个近似线性的片；微分降秩时，则必须证明像的 \(m\) 维面积确实消失。后一步是连接线性公式与一般公式的另一半。

\begin{lemma}\label{g3:lem:rank-deficient-image-null}
	设 \(1\le m\le n\)，\(f:\mathbb R^m\to\mathbb R^n\) 为 Lipschitz 映射，令
	\[
		C=\{x:f\text{ 在 }x\text{ 可微，且 }\operatorname{rank}\mathrm Df(x)<m\}.
	\]
	则 \(\mathcal H^m(f(C))=0\)。
\end{lemma}
\begin{proof}
	记 \(L=\operatorname{Lip}f\)，先固定 \(R>0\)，只处理 \(C_R=C\cap B_R(0)\)。取 \(0<\varepsilon<1\)。对每个 \(x\in C_R\)，可微性给出任意小的半径 \(r>0\)，使
	\[
		|f(x+h)-f(x)-T_xh|\le\varepsilon|h|\quad(|h|\le5r),\qquad T_x=\mathrm Df(x).
	\]
	还要求 \(r<1/5\)，以及 \(r\) 小于一个稍后趋零的统一上界。由于 \(\lVert T_x\rVert\le L\)，\(f(B_{5r}(x))\) 包含在仿射平面 \(f(x)+\operatorname{im}T_x\) 中半径 \(5Lr\) 的圆盘的 \(5\varepsilon r\) 邻域内。

	设 \(k=\operatorname{rank}T_x\le m-1\)。在这个 \(k\) 维平面内取边长 \(\varepsilon r\) 的立方体网格，半径 \(5Lr\) 的圆盘只遇到至多
	\[
		C_{m,n}(1+L)^k\varepsilon^{-k}
	\]
	个网格块。每个网格块的 \(5\varepsilon r\) 邻域可放入一个直径至多 \(C_{m,n}\varepsilon r\) 的 \(n\) 维球中。因此该像集有一个覆盖，其 \(m\) 维代价至多为
	\[
		C_{m,n}(1+L)^k\varepsilon^{m-k}r^m
		\le C_{m,n}(1+L)^m\varepsilon r^m.
	\]
	当 \(k=0\) 时只需一个球，同一估计仍成立。这正是降秩的作用：至少有一个方向提供额外的 \(\varepsilon\) 因子。

	对所有满足上述条件、以 \(C_R\) 中的点为中心的闭球应用\cref{g3:lem:five-ball-covering}，选出两两不交的闭球 \(\overline B_{r_i}(x_i)\)，使 \(C_R\subset\bigcup_i\overline B_{5r_i}(x_i)\)。前面的可微性估计同样控制这些五倍闭球的像。这些原球均位于 \(B_{R+1}(0)\)，所以
	\[
		\sum_i\omega_m r_i^m\le\mathcal L^m(B_{R+1}(0)).
	\]
	给定任意 \(\delta>0\)，预先把半径上界取得足够小，使前述覆盖球的直径都不超过 \(\delta\)。将各球像的细长盒覆盖合并，便得
	\[
		\mathcal H^m_\delta(f(C_R))
		\le C_{m,n,L}\varepsilon(R+1)^m.
	\]
	常数与 \(\delta,\varepsilon\) 无关。先令 \(\delta\downarrow0\)，再令 \(\varepsilon\downarrow0\)，得到 \(\mathcal H^m(f(C_R))=0\)。最后令整数 \(R\uparrow\infty\)，用可数次可加性完成证明。
\end{proof}

\begin{theorem}\label{g3:thm:injective-lipschitz-area}
	设 \(1\le m\le n\)，\(f:\mathbb R^m\to\mathbb R^n\) 为 Lipschitz 映射，\(A\subset\mathbb R^m\) 为 Lebesgue 可测集。若 \(f|_A\) 单射，则
	\begin{equation}\label{g3:eq:injective-lipschitz-area}
		\mathcal H^m(f(A))=\int_AJ_mf(x)\,\mathrm d\mathcal L^m(x).
	\end{equation}
\end{theorem}
\begin{proof}
	设 \(D\) 为 \(f\) 的可微集。由上一节的可微性定理，\(A\setminus D\) 为零测集，其像也为 \(\mathcal H^m\) 零测集。由\cref{g3:lem:rank-deficient-image-null}，\(A\cap C\) 的像同样为零测集，而 \(J_mf\) 在这里等于零。因此只需处理满秩部分。

	对整数 \(s\ge1\)，令
	\[
		A_s=\{x\in A\cap D:\ |\mathrm Df(x)v|\ge s^{-1}|v|\text{ 对所有 }v\in\mathbb R^m\}.
	\]
	这些集可测，随 \(s\) 递增，且并集恰为 \(A\) 中的满秩可微点。固定 \(s\)，再取 \(0<\varepsilon<(4s)^{-1}\)。用\cref{g3:thm:lipschitz-linearization-pieces}把 \(A_s\) 分成可数个互不交可测片 \(E_i\)。每片有固定矩阵 \(T_i\)，满足
	\[
		\lVert\mathrm Df(x)-T_i\rVert\le\varepsilon,\qquad
		|f(y)-f(x)-T_i(y-x)|\le\varepsilon|y-x|\quad(x,y\in E_i).
	\]
	删除空片。由 \(E_i\subset A_s\)，\(|T_iv|\ge(s^{-1}-\varepsilon)|v|\ge(2s)^{-1}|v|\)。记 \(\rho=2s\varepsilon<1/2\)。\Cref{g3:lem:near-linear-area-comparison}给出
	\[
		(1-\rho)^mJ_m(T_i)\mathcal L^m(E_i)
		\le\mathcal H^m(f(E_i))
		\le(1+\rho)^mJ_m(T_i)\mathcal L^m(E_i),
	\]
	并且对 \(x\in E_i\)，\((1-\rho)^mJ_m(T_i)\le J_mf(x)\le(1+\rho)^mJ_m(T_i)\)。合并得到
	\[
		\left(\frac{1-\rho}{1+\rho}\right)^m\int_{E_i}J_mf\,\mathrm d\mathcal L^m
		\le\mathcal H^m(f(E_i))
		\le\left(\frac{1+\rho}{1-\rho}\right)^m\int_{E_i}J_mf\,\mathrm d\mathcal L^m.
	\]
	由\cref{g3:prop:area-multiplicity-measurable}，这些像均可测；由于 \(f|_A\) 单射，它们还两两不交。因此可对 \(i\) 求和，以 \(A_s\) 代替上式中的 \(E_i\)。再令 \(\varepsilon\downarrow0\)，两侧常数趋于一，得到 \(A_s\) 上的等式。此论证同样适用于积分为无穷的情形，因为下界已迫使像测度为无穷。

	最后令 \(s\uparrow\infty\)，对参数域上的积分及递增像集的测度分别使用单调收敛，再加回前述两个像零测的部分，得到式\eqref{g3:eq:injective-lipschitz-area}。
\end{proof}

\subsection{非单射情形的分片求和}
\label{g3:subsec:2003:04}
现在允许不同的参数点具有相同像。各单射片的面积仍能分别计算；它们在目标中相遇时，示性函数相加，恰好产生重数。

\begin{theorem}[面积公式]\label{g3:thm:lipschitz-area-formula}
	设 \(1\le m\le n\)，\(f:\mathbb R^m\to\mathbb R^n\) 为 Lipschitz 映射，\(A\subset\mathbb R^m\) 为 Lebesgue 可测集。则
	\begin{equation}\label{g3:eq:lipschitz-area-formula}
		\int_AJ_mf\,\mathrm d\mathcal L^m
		=\int_{\mathbb R^n}N(f,A,y)\,\mathrm d\mathcal H^m(y).
	\end{equation}
	更一般地，对非负 Lebesgue 可测函数 \(a:A\to[0,\infty]\)，
	\begin{equation}\label{g3:eq:weighted-lipschitz-area-formula}
		\int_Aa(x)J_mf(x)\,\mathrm d\mathcal L^m(x)
		=\int_{\mathbb R^n}\left(\sum_{x\in A\cap f^{-1}(y)}a(x)\right)\mathrm d\mathcal H^m(y).
	\end{equation}
	两边允许取 \(+\infty\)，非负积分中的 \(0\cdot\infty\) 取零。
\end{theorem}
\begin{proof}
	把 \(A\) 写成 \(Z\cup C_A\cup\bigcup_iA_i\)：其中 \(Z\) 是不可微点，\(C_A\) 是可微但降秩的点，\(A_i\) 是满秩部分的互不交可测双 Lipschitz 片。这一分解由\cref{g3:cor:full-rank-bilipschitz-pieces}给出；若原覆盖有重叠，依次从第 \(i\) 片扣除前面的片即可。前述证明表明
	\[
		\mathcal H^m(f(Z\cup C_A))=0,
		\qquad \int_{Z\cup C_A}J_mf\,\mathrm d\mathcal L^m=0.
	\]
	每个 \(A_i\) 上 \(f\) 单射。对 \(y\notin f(Z\cup C_A)\)，每个原像点落在唯一一片中，所以
	\[
		N(f,A,y)=\sum_i\mathbf1_{f(A_i)}(y).
	\]
	因此，由可测性、单射面积公式和单调收敛，
	\[
		\int_{\mathbb R^n}N(f,A,y)\,\mathrm d\mathcal H^m(y)
		=\sum_i\mathcal H^m(f(A_i))
		=\sum_i\int_{A_i}J_mf\,\mathrm d\mathcal L^m
		=\int_AJ_mf\,\mathrm d\mathcal L^m.
	\]
	这证明无权公式。对简单函数 \(a\)，在其各个可测水平集上使用该公式并相加，就得到带权公式。一般非负 \(a\) 取简单函数 \(a_j\uparrow a\)，由\cref{g3:prop:area-multiplicity-measurable}的最后一步，纤维上的和也递增到所需函数。两侧再次使用单调收敛即可。
\end{proof}

当函数有符号或取复值时，只要 \(\int_A|a|J_mf\,\mathrm d\mathcal L^m<\infty\)，就可分别对正负部或实虚部应用公式。带绝对值的公式同时保证目标中纤维和几乎处处绝对收敛。

\begin{corollary}\label{g3:cor:area-formula-subsets}
	设 \(1\le m\le n\)，\(E\subset\mathbb R^m\) 为 Lebesgue 可测集，\(f:E\to\mathbb R^n\) 为 Lipschitz 映射。在 \(E\) 的密度一点，以任意全空间 Lipschitz 延拓的微分定义 \(\mathrm D_Ef\)，在其余零测例外点任意补零。则 \(\mathrm D_Ef\) 在 \(E\) 上几乎处处存在并且与延拓无关。以 \(J_mf=\sqrt{\det((\mathrm D_Ef)^{\mathsf T}\mathrm D_Ef)}\) 记面积因子，式\eqref{g3:eq:lipschitz-area-formula}和\eqref{g3:eq:weighted-lipschitz-area-formula}对任意可测 \(A\subset E\) 及非负可测权仍成立。
\end{corollary}
\begin{proof}
	若 \(E\) 为空，结论直接成立。否则由\cref{g3:prop:lipschitz-extension-density}，对每个坐标作全空间 Lipschitz 延拓即可得到 \(F:\mathbb R^m\to\mathbb R^n\)。该命题还证明，在 \(E\) 的密度一点，只要一个延拓可微，所有延拓就有同一微分；密度定理和可微性定理说明这排除的只是 \(E\) 中的零测集。对 \(F\) 和参数集 \(A\) 应用面积公式，其在 \(A\) 上的值及微分恰是所需的值与微分。
\end{proof}

上述陈述也解释了零维端点：\(\mathbb R^0\) 是一个点，令 \(J_0=1\)，源、靶所用 \(\mathcal H^0\) 均为计数测度，则一个零维参数片的带权面积公式就是有限单点和的恒等式。

\begin{corollary}\label{g3:cor:bilipschitz-change-variables}
	设 \(m\ge1\)，\(E\subset\mathbb R^m\) 为 Lebesgue 可测集，\(f:E\to\mathbb R^m\) 为到其像的双 Lipschitz 映射。则 \(f(E)\) 可测，双方的零测集在 \(f\) 与 \(f^{-1}\) 下仍为零测集；对几乎处处的 \(x\in E\)，双方的微分均存在，且
	\begin{equation}\label{g3:eq:bilipschitz-inverse-differential}
		\mathrm D_{f(E)}f^{-1}(f(x))\,\mathrm D_Ef(x)=I_m.
	\end{equation}
	特别地，\(J_mf(x)=|\det\mathrm D_Ef(x)|>0\) 几乎处处。对可测 \(A\subset E\) 及非负可测 \(h:f(E)\to[0,\infty]\)，有
	\begin{equation}\label{g3:eq:bilipschitz-change-variables}
		\int_{f(A)}h(y)\,\mathrm d\mathcal L^m(y)
		=\int_Ah(f(x))\,|\det\mathrm D_Ef(x)|\,\mathrm d\mathcal L^m(x).
	\end{equation}
\end{corollary}
\begin{proof}
	像的可测性由延拓及\cref{g3:prop:area-multiplicity-measurable}得到。零测集的像为零测，直接来自双向 Lipschitz 控制与 \(\mathcal H^m=\mathcal L^m\)；同理，任意可测集的原像也可测。于是 \(h\circ f\) 是可测函数。

	分别延拓 \(f\) 与 \(f^{-1}\) 为全空间 Lipschitz 映射 \(F,G\)。除去源中的非密度点、\(F\) 的不可微点，以及目标中非密度点和 \(G\) 不可微点的原像后，仅损失一个零测集。在其余 \(x\) 上，\(G\circ F\) 可微，并在 \(E\) 上等于恒等映射。由密度点上延拓微分的唯一性，\(\mathrm D(G\circ F)(x)=I_m\)。通常的可微映射链式法则给出式\eqref{g3:eq:bilipschitz-inverse-differential}，故两个方阵均可逆。最后对单射映射 \(f\) 使用带权面积公式，并取权 \(h\circ f\)，即可得到换元公式。
\end{proof}

\subsection{几何面积与参数面积的区别}
\label{g3:subsec:2003:05}
面积公式给出两个不同量的精确关系：
\[
	\underbrace{\mathcal H^m(f(A))}_{\text{像集的几何面积}}
	\le
	\underbrace{\int_AJ_mf\,\mathrm d\mathcal L^m}_{\text{参数面积}}
	=\int_{f(A)}N(f,A,y)\,\mathrm d\mathcal H^m(y).
\]
当参数面积有限时，两者相等当且仅当 \(N(f,A,y)=1\) 对 \(\mathcal H^m\) 几乎处处的 \(y\in f(A)\) 成立。这由 \(N\ge1\) 及非负函数 \(N-1\) 的积分为零直接得到。若两边均为无穷，相等本身不能推出重数为一。

\begin{example}\label{g3:exm:area-fold-and-repetition}
	取 \(f:[-1,1]\to\mathbb R\)、\(f(t)=|t|\)。除 \(t=0\) 外，\(J_1f=1\)，故参数长度为 \(2\)。像是 \([0,1]\)，几何长度为 \(1\)，而每个 \(y\in(0,1)\) 有两个原像。这里是折叠造成重复。另取 \(\gamma:[0,4\pi]\to\mathbb R^2\)、\(\gamma(t)=(\cos t,\sin t)\)。它处处有速度一，参数长度为 \(4\pi\)，像圆的长度为 \(2\pi\)；除端点所对应的一个像点外，重数为二。这里即使每处微分满秩，整体重复仍然存在。
\end{example}

\begin{proposition}\label{g3:prop:lipschitz-graph-area}
	设 \(m,q\ge1\)，\(E\subset\mathbb R^m\) 为 Lebesgue 可测集，\(u:E\to\mathbb R^q\) 为 Lipschitz 映射，\(G(x)=(x,u(x))\)。则对每个可测 \(A\subset E\)，
	\begin{equation}\label{g3:eq:lipschitz-graph-area}
		\mathcal H^m(G(A))
		=\int_A\sqrt{\det\bigl(I_m+(\mathrm D_Eu)^{\mathsf T}\mathrm D_Eu\bigr)}\,\mathrm d\mathcal L^m.
	\end{equation}
	当 \(q=1\) 时，被积函数为 \(\sqrt{1+|\mathrm D_Eu|^2}\)。
\end{proposition}
\begin{proof}
	图映射单射，逆为第一分量投影。几乎处处有
	\[
		\mathrm D_EG=\begin{pmatrix}I_m\\\mathrm D_Eu\end{pmatrix},\qquad
		(\mathrm D_EG)^{\mathsf T}\mathrm D_EG=I_m+(\mathrm D_Eu)^{\mathsf T}\mathrm D_Eu.
	\]
	代入单射面积公式即可。若 \(q=1\)，矩阵 \((\mathrm D_Eu)^{\mathsf T}\mathrm D_Eu\) 在梯度方向的特征值为 \(|\mathrm D_Eu|^2\)，在其正交补上为零，故行列式为 \(1+|\mathrm D_Eu|^2\)。
\end{proof}

光滑参数曲面的面积元 \(\sqrt{\det(\langle\partial_iF,\partial_jF\rangle)}\,\mathrm dx\) 是同一公式的特例；图公式则把第十九章的光滑图面积延伸到 Lipschitz 图。参数域换成双 Lipschitz 坐标也不改变面积积分：若 \(\phi:E'\to E\subset\mathbb R^m\) 为双 Lipschitz 映射，\(F:E\to\mathbb R^n\) 为 Lipschitz 映射，则在双方密度点用延拓和链式法则，有
\[
	J_m(F\circ\phi)(x)
	=J_mF(\phi(x))\,|\det\mathrm D_{E'}\phi(x)|.
\]
这里的线性恒等式来自
\(\det(B^{\mathsf T}A^{\mathsf T}AB)=(\det B)^2\det(A^{\mathsf T}A)\)。异常点在双 Lipschitz 换元下仍为零测，再用式\eqref{g3:eq:bilipschitz-change-variables}即可比较两个参数积分。
\ProseParagraphBreak
因而“重参数化”与“重复覆盖”有不同的面积后果：前者以双射改变坐标，后者使像点的原像数增加。接下来讨论切片时，面积公式会用于每一张低维单射参数片；讨论可整流集时，则用它把可数参数片上的微分与像集上的几何测度联系起来。
% 纲要条目排为编号小节标题，并进入目录与书签；说明段仍保留于纲要注释中。

% !TeX root = ../../Book3.tex
% 纲要标题：### 20.4 余面积公式
% 纲要位置：几何学计划纲要.md，第 2948 行。
% 正文状态：初稿；尚未正式审核。

% 纲要内容开始
% 
% * 线性投影模型
% * 切片的测度与 Jacobian
% * 局部分解和积分交换
% * 光滑水平集及非光滑切片
% 
% 纲要内容结束

\section{余面积公式}
\label{g3:sec:2004}

面积公式把参数域上的积分送到像上；当源空间的维数更大时，像点通常对应一整片纤维，单纯计数已经不够。余面积公式把每根纤维的低维面积作为内层积分，再对像空间积分。正交投影时这就是 Fubini 定理，非线性映射则需要补上横向伸缩率，并说明秩亏点和不可微点如何参与切片。

\ProseParagraphBreak

本节固定整数 \(m\ge n\ge1\)，记 \(k=m-n\)。\(\mathcal L^j\) 表示 \(j\) 维 Lebesgue 测度，\(\mathcal H^j\) 使用\cref{g3:thm:hausdorff-lebesgue}的标准归一化；特别地，它在 \(j\) 维仿射平面上就是 Euclid 体积。\(\mathcal H^0\) 为计数测度，\(\mathbb R^0\) 由一个点组成，赋予质量一。非负积分允许为 \(+\infty\)，并约定 \(0\cdot\infty=0\)。

\subsection{线性投影模型}
\label{g3:subsec:2004:01}
把 \(x\in\mathbb R^m\) 写成 \(x=(y,z)\in\mathbb R^n\times\mathbb R^k\)。投影 \(p(y,z)=y\) 的纤维是 \(\{y\}\times\mathbb R^k\)，故对任意非负 Lebesgue 可测函数 \(g\)，非负积分的 Fubini 定理给出
\[
	\int_{\mathbb R^m}g(x)\,\mathrm dx
	=\int_{\mathbb R^n}\left(\int_{p^{-1}(y)}g\,\mathrm d\mathcal H^k\right)\mathrm dy.
\]
内层测度的维数是剩下的 \(k\)，外层测度的维数是输出的 \(n\)。对于一般线性映射，应当测量它在核的正交补上把 \(n\) 维体积放大了多少。

\begin{definition}\label{g3:def:coarea-jacobian}
	对线性映射 \(A:\mathbb R^m\to\mathbb R^n\)，其中 \(m\ge n\ge1\)，定义
	\[
		J_n(A)=\sqrt{\det(AA^{\mathsf T})}.
	\]
	若 \(f:\mathbb R^m\to\mathbb R^n\) 为 Lipschitz 映射，则在可微点令 \(J_nf(x)=J_n(\mathrm Df(x))\)，在其余点令 \(J_nf(x)=0\)。这个函数称为 \(f\) 的 \(n\) 维\term[yu mianji Jacobian]{余面积 Jacobian}[coarea Jacobian]。
\end{definition}
\symidx{\(J_n f\)：目标维数为 \(n\) 的余面积 Jacobian}

按行列式乘积的子式展开，\(J_n(A)^2\) 等于 \(A\) 的所有 \(n\) 阶子式的平方和。因此 \(J_n(A)>0\) 恰当 \(\operatorname{rank}A=n\)；当 \(n=1\) 时，\(J_1f=|\nabla f|\)，当 \(n=m\) 时，\(J_mf=|\det\mathrm Df|\)。若 \(f\) 的 Lipschitz 常数为 \(L\)，则几乎处处有 \(0\le J_nf\le L^n\)。这些量由可测的导数矩阵给出，因而可测。

\begin{proposition}\label{g3:prop:linear-coarea}
	设 \(A:\mathbb R^m\to\mathbb R^n\) 为线性映射，\(m\ge n\ge1\)、\(k=m-n\)，\(g:\mathbb R^m\to[0,+\infty]\) 为 Lebesgue 可测函数。则纤维积分几乎处处有定义且为 Lebesgue 可测函数，并有
	\[
		J_n(A)\int_{\mathbb R^m}g\,\mathrm dx
		=\int_{\mathbb R^n}\left(\int_{A^{-1}(y)}g\,\mathrm d\mathcal H^k\right)\mathrm dy.
	\]
\end{proposition}
\begin{proof}
	先设 \(A\) 满秩。在 \((\ker A)^\perp\) 与 \(\ker A\) 中分别选正交基，写成矩阵 \(V\) 与 \(U\)，则 \(Q=(V,U)\) 正交，且 \(AQ=(B,0)\)，其中 \(B=AV\) 可逆。于是 \(J_n(A)=|\det B|\)。作正交换元 \(x=Vs+Uz\)，再作 \(y=Bs\)，得到
	\[
		J_n(A)\int g(x)\,\mathrm dx
		=\int_{\mathbb R^n}\int_{\mathbb R^k}
		g(VB^{-1}y+Uz)\,\mathrm dz\,\mathrm dy.
	\]
	对固定 \(y\)，映射 \(z\mapsto VB^{-1}y+Uz\) 是到纤维的等距参数化，所以内层正是所需 \(\mathcal H^k\) 积分。Fubini 定理也给出切片可测性和外层可测性。
	若 \(\operatorname{rank}A<n\)，则 \(J_n(A)=0\)，且纤维在 \(\operatorname{im}A\) 之外为空；这个真线性子空间的 \(\mathcal L^n\) 测度为零。因此右端在几乎处处为零，对例外集上的切片积分任取零的代表，等式两端均为零。
\end{proof}

下面的矩阵恒等式把横向伸缩与纤维内的面积联系起来，稍后会逐点应用于非线性映射。

\begin{lemma}\label{g3:lem:coarea-linear-block-identity}
	设 \(m\ge n\ge1\)、\(k=m-n\)，\(A:\mathbb R^m\to\mathbb R^n\) 满秩，\(P:\mathbb R^m\to\mathbb R^k\) 为线性映射，并且
	\[
		T=\begin{pmatrix}A\\P\end{pmatrix}:\mathbb R^m\longrightarrow\mathbb R^n\times\mathbb R^k
	\]
	可逆。令 \(\iota(z)=(0,z)\)、\(B=T^{-1}\iota\)。则
	\begin{equation}\label{g3:eq:coarea-linear-block-identity}
		J_n(A)=|\det T|\sqrt{\det(B^{\mathsf T}B)}.
	\end{equation}
	当 \(k=0\) 时，右端的空行列式取一。
\end{lemma}
\begin{proof}
	沿用上一证明的正交矩阵 \((V,U)\)。因为
	\[
		T(V,U)=\begin{pmatrix}AV&0\\PV&PU\end{pmatrix},
	\]
	所以 \(PU\) 可逆，且 \(|\det T|=|\det(AV)|\,|\det(PU)|\)。另一方面，\(AB=0\)、\(PB=I\) 唯一确定 \(B=U(PU)^{-1}\)，故
	\[
		\sqrt{\det(B^{\mathsf T}B)}=|\det(PU)|^{-1}.
	\]
	再用 \(|\det(AV)|=J_n(A)\) 即得结论。\(k=0\) 时公式直接退化为方阵行列式的绝对值。
\end{proof}

\subsection{切片的测度与 Jacobian}
\label{g3:subsec:2004:02}
非线性切片首先遇到一个测度问题：一个 \(m\) 维零测集可能含有整张 \(k\) 维曲面，不能在每一根纤维上直接删去。所需结论是，它在几乎所有纤维上具有零 \(k\) 维测度。先用一个不要求切片函数已知可测的覆盖估计来证明这一点。

\ProseParagraphBreak

对于任意非负函数 \(h\)，记 \(\int^*h\) 为其 Lebesgue 上积分，即所有满足 \(v\ge h\) 的非负可测函数 \(v\) 的积分下确界。若 \(h\) 可测，它就是普通积分。以下记 \(c_j=\omega_j/2^j\)，其中 \(c_0=1\)，并以 \(\mathcal H^j_\delta\) 表示由直径小于 \(\delta\) 的覆盖定义的 Hausdorff 近似外测度。

\begin{lemma}\label{g3:lem:coarea-covering-inequality}
	设 \(f:\mathbb R^m\to\mathbb R^n\) 的 Lipschitz 常数为 \(L\)，\(m\ge n\ge1\)、\(k=m-n\)。以下 \(C_{m,n}\) 为只依赖维数的常数，\(\int^*\) 表示非负函数的上积分。对任意集合 \(E\subset\mathbb R^m\)，有
	\begin{equation}\label{g3:eq:coarea-covering-inequality}
		\int_{\mathbb R^n}^{*}\mathcal H^k(E\cap f^{-1}(y))\,\mathrm dy
		\le C_{m,n}L^n\mathcal H^m(E).
	\end{equation}
	特别地，若 \(\mathcal L^m(E)=0\)，则 \(\mathcal H^k(E\cap f^{-1}(y))=0\) 对几乎所有 \(y\) 成立。
\end{lemma}
\begin{proof}
	先设 \(L>0\)、\(\mathcal H^m(E)<\infty\)。取覆盖 \(E\subset\bigcup_iE_i\)，其中 \(d_i=\operatorname{diam}E_i<\delta\)，并删去空的 \(E_i\)。在各 \(E_i\) 中取一点 \(x_i\)。若 \(E_i\) 与纤维 \(f^{-1}(y)\) 相交，则 \(y\) 属于闭球 \(B_i=\overline B(f(x_i),Ld_i)\)。因此
	\[
		\mathcal H^k_\delta(E\cap f^{-1}(y))
		\le c_k\sum_i d_i^k\mathbf1_{B_i}(y)=:v_\delta(y),
	\]
	其中 \(d_i^0\) 一律取一。即使 \(d_i=0\)，其对应的目标球也只有一个点，不影响下面的积分。逐项积分给出
	\[
		\int v_\delta\le c_k\omega_nL^n\sum_i d_i^{k+n}
		=\frac{c_k\omega_n}{c_m}L^n\sum_i c_md_i^m.
	\]
	选 \(\delta_j\downarrow0\)，并对每个 \(j\) 选覆盖，使末尾的和不超过 \(\mathcal H^m(E)+1/j\)。对相应的可测函数 \(v_j\)，有
	\[
		\mathcal H^k(E\cap f^{-1}(y))
		\le\liminf_jv_j(y).
	\]
	右端可测，Fatou 引理于是证明式\eqref{g3:eq:coarea-covering-inequality}，常数可取 \(c_k\omega_n/c_m\)。若 \(\mathcal H^m(E)=\infty\)，不等式直接成立。若 \(L=0\)，映射恒定，所有非空纤维集中在一个目标点上，上积分为零，也得到结论。
	最后使用 \(\mathcal H^m=\mathcal L^m\)。当右端为零时，任意 \(a>0\) 所对应集合 \(\{y:\mathcal H^k(E\cap f^{-1}(y))>a\}\) 的 Lebesgue 外测度为零：用积分任意小的可测上界和 Chebyshev 不等式即可。对 \(a=1/j\) 取并，得到最后的断言。
\end{proof}

这个估计尚未确定正确的 Jacobian，但它已经允许在切片中删去源空间的零测集。利用紧集逼近，就能解决可测性。

\begin{lemma}\label{g3:lem:coarea-slice-measurability}
	设 \(f:\mathbb R^m\to\mathbb R^n\) 为 Lipschitz 映射，\(m\ge n\ge1\)、\(k=m-n\)。若 \(E\) 为 Lebesgue 可测集，则 \(E\cap f^{-1}(y)\) 对几乎所有 \(y\) 为 \(\mathcal H^k\) 可测集，且
	\[
		y\longmapsto\mathcal H^k(E\cap f^{-1}(y))
	\]
	在取几乎处处代表后 Lebesgue 可测。更一般地，对任意非负 Lebesgue 可测 \(g\)，它在几乎所有纤维上为 \(\mathcal H^k\) 可测函数，且 \(y\mapsto\int_{f^{-1}(y)}g\,\mathrm d\mathcal H^k\) Lebesgue 可测。
\end{lemma}
\begin{proof}
	先取紧集 \(K\)。令 \(h_\delta(y)\) 为所有覆盖 \(K\cap f^{-1}(y)\) 的有限族开集 \(U_1,\ldots,U_N\) 所对应 \(\sum_i c_k(\operatorname{diam}U_i)^k\) 的下确界，要求各直径小于 \(\delta\)；空纤维允许空覆盖。固定一族这样的开集，它在 \(y\) 附近仍覆盖纤维：否则有 \(y_j\to y\) 和 \(x_j\in K\cap f^{-1}(y_j)\) 落在全部 \(U_i\) 外，紧性给出极限点 \(x\in K\cap f^{-1}(y)\)，矛盾。因此 \(\{h_\delta<a\}\) 为开集，\(h_\delta\) 上半连续。
	我们还需核对 \(\lim_{\delta\downarrow0}h_\delta(y)=\mathcal H^k(K\cap f^{-1}(y))\)。一边由 \(\mathcal H^k_\delta\le h_\delta\) 得到。另一边，把一个直径小于 \(\delta/2\) 的可数覆盖中的集合略微加厚为开集，使新直径仍小于 \(\delta\)，且总费用增加至多任给的 \(\eta>0\)；逐项选择增加量小于 \(\eta2^{-i}\) 即可。紧性允许选有限子覆盖，费用只会减小。因此 \(h_\delta\le\mathcal H^k_{\delta/2}+\eta\)。\(k=0\) 时，加厚不改变每项费用一，同一论证适用。由此得到所需极限和紧集情形的可测性。
	对一般 Lebesgue 可测 \(E\)，内正则性与有界区域穷尽给出递增紧集 \(K_j\subset E\)，使 \(N=E\setminus\bigcup_jK_j\) 为 \(\mathcal L^m\) 零测集。由\cref{g3:lem:coarea-covering-inequality}，\(N\cap f^{-1}(y)\) 对几乎所有 \(y\) 为 \(\mathcal H^k\) 零测集。Hausdorff 测度的完备性说明这些 \(E\cap f^{-1}(y)\) 可测，而且
	\[
		\mathcal H^k(E\cap f^{-1}(y))
		=\lim_j\mathcal H^k(K_j\cap f^{-1}(y)).
	\]
	这证明集合版本。最后，非负 Lebesgue 可测函数在一个 \(\mathcal L^m\) 零测集之外有 Borel 代表；它可由非负简单函数单调逼近。把刚才的零测切片结论和集合版本应用于这些函数所涉及的可数个集合，再用单调收敛，即得函数版本。
\end{proof}

下面处理导数存在但秩不足的点。关键是它们的小球像在至少一个目标方向上很薄。源球本身仍是 \(m\) 维的；薄的是目标中的覆盖范围，这正好提供积分中消失的因子。

\begin{lemma}\label{g3:lem:coarea-critical-slices}
	设 \(f:\mathbb R^m\to\mathbb R^n\) 为 Lipschitz 映射，\(m\ge n\ge1\)、\(k=m-n\)，并令
	\[
		C=\{x:f\text{ 在 }x\text{ 可微且 }\operatorname{rank}\mathrm Df(x)<n\}.
	\]
	则 \(\mathcal H^k(C\cap f^{-1}(y))=0\) 对 \(\mathcal L^n\) 几乎所有 \(y\) 成立。
\end{lemma}
\begin{proof}
	固定 \(R>0\)，只处理 \(C_R=C\cap B(0,R)\)。取 \(0<\varepsilon<1\)、\(0<\delta<1\)，记 \(L\) 为 \(f\) 的 Lipschitz 常数。对每个 \(x\in C_R\)，可微性允许任取充分小的 \(r>0\)，使 \(10r<\delta\)，并对 \(|h|\le5r\) 有
	\[
		|f(x+h)-f(x)-\mathrm Df(x)h|\le\varepsilon|h|.
	\]
	若 \(q=\operatorname{rank}\mathrm Df(x)\le n-1\)，在目标中选前 \(q\) 个坐标张成该导数的像，则 \(f(\overline B(x,5r))\) 包含于一个矩形盒 \(Q_{x,r}\)：其前 \(q\) 个边长不超过 \(10(L+1)r\)，其余边长不超过 \(10\varepsilon r\)。因而
	\begin{equation}\label{g3:eq:coarea-critical-thin-box}
		\mathcal L^n(Q_{x,r})
		\le10^n(L+1)^n\varepsilon r^n.
	\end{equation}
	使用\cref{g3:lem:five-ball-covering}，从上述闭球 \(\overline B(x,r)\) 的族中选出两两不交的 \(\overline B(x_i,r_i)\)，使 \(C_R\subset\bigcup_i\overline B(x_i,5r_i)\)。这些球的中心在 \(B(0,R)\)、半径小于一，因此
	\[
		\sum_i\omega_mr_i^m\le\mathcal L^m(B(0,R+1)).
	\]
	每个纤维与第 \(i\) 个扩大球的交的直径小于 \(\delta\)，且只有 \(y\in Q_{x_i,r_i}\) 时该交才可能非空。因此可测函数
	\[
		v_\delta(y)=c_k\sum_i(10r_i)^k\mathbf1_{Q_{x_i,r_i}}(y)
	\]
	满足 \(\mathcal H^k_\delta(C_R\cap f^{-1}(y))\le v_\delta(y)\)。由式\eqref{g3:eq:coarea-critical-thin-box}和上面的球体积和，
	\[
		\int v_\delta\le C(m,n,L)\varepsilon(R+1)^m.
	\]
	现在令 \(\delta\downarrow0\)，如\cref{g3:lem:coarea-covering-inequality}的证明那样对一列可测上界取下极限并使用 Fatou 引理，得到
	\[
		\int^*\mathcal H^k(C_R\cap f^{-1}(y))\,\mathrm dy
		\le C(m,n,L)\varepsilon(R+1)^m.
	\]
	因 \(\varepsilon\) 任意，上积分为零。最后令 \(R\) 取正整数并取可数并，就得到结论。此证明也适用于 \(k=0\)，此时每个非空覆盖片的费用为一。
\end{proof}

\subsection{局部分解和积分交换}
\label{g3:subsec:2004:03}
剩下的是满秩可微点。在这些点附近，选出导数的一个非零 \(n\) 阶子式，并把其余源坐标补到输出后面，就得到方阵映射。它未必在整个邻域内单射，但\cref{g3:cor:full-rank-bilipschitz-pieces}使它在可数片上双 Lipschitz。这已经足够逐片换元。

\begin{lemma}\label{g3:lem:coarea-regular-pieces}
	设 \(f:\mathbb R^m\to\mathbb R^n\) 为 Lipschitz 映射，\(m\ge n\ge1\)、\(k=m-n\)。除去一个 \(\mathcal L^m\) 零测集后，其满秩可微点集合可分为可数个互不相交的 Borel 集 \(E_j\)，使每个 \(E_j\) 对应一个坐标投影 \(P_j:\mathbb R^m\to\mathbb R^k\)，且
	\[
		G_j=(f,P_j):E_j\longrightarrow B_j:=G_j(E_j)
	\]
	为双 Lipschitz 双射。还可要求 \(E_j\)、\(B_j\) 均为可数个紧集的并。
\end{lemma}
\begin{proof}
	在满秩点选 \(\mathrm Df\) 的 \(n\) 个线性无关列，令 \(P\) 由其余 \(k\) 个源坐标组成，则 \(\mathrm D(f,P)\) 可逆。只有有限种列选择。对每个全空间 Lipschitz 映射 \(G=(f,P)\) 应用\cref{g3:cor:full-rank-bilipschitz-pieces}，再与所考虑的满秩点集合相交，便得到可数个双 Lipschitz Borel 片，覆盖这一集合。
	在每个片与有界球的交中用内正则性选紧子集，留下的部分总测度为零。把所得全部紧集列成 \(K_1,K_2,\ldots\)，置 \(E_j=K_j\setminus\bigcup_{i<j}K_i\)，则各片互不相交，且各自仍继承所属 \(G\) 的双 Lipschitz 性。因为 \(E_j\) 是紧集 \(K_j\) 内的相对开集，它是可数个紧集的并；其连续像 \(B_j\) 也有这一性质，特别地都是 Borel 集。
\end{proof}

先在一个这样的片上证明公式。逆映射的纤维内导数将产生 \(k\) 维面积因子，而全维换元产生一个行列式；式\eqref{g3:eq:coarea-linear-block-identity}正好识别二者的乘积。

\begin{lemma}\label{g3:lem:coarea-on-bilipschitz-piece}
	设 \(f:\mathbb R^m\to\mathbb R^n\) 为 Lipschitz 映射，\(m\ge n\ge1\)，\(k=m-n\)。设 \(E\subset\mathbb R^m\) 为可数紧集之并，存在坐标投影 \(P:\mathbb R^m\to\mathbb R^k\)，使 \((f,P)|_E\) 为到其像的双 Lipschitz 映射；这包括\cref{g3:lem:coarea-regular-pieces}中的各片。对任意非负 Lebesgue 可测函数 \(g:\mathbb R^m\to[0,\infty]\)，有
	\[
		\int_EgJ_nf\,\mathrm dx
		=\int_{\mathbb R^n}\left(\int_{E\cap f^{-1}(y)}g\,\mathrm d\mathcal H^k\right)\mathrm dy.
	\]
\end{lemma}
\begin{proof}
	记 \(G=(f,P)\)、\(B=G(E)\)，逆映射为 \(h:B\to E\)。用\cref{g3:prop:lipschitz-extension-density}逐坐标把 \(h\) 延拓为全空间 Lipschitz 映射 \(\widetilde h\)。由\cref{g3:cor:bilipschitz-change-variables}，对几乎所有 \(w=(y,z)\in B\)，\(x=h(w)\) 是 \(E\) 的密度点，并有
	\[
		\mathrm D\widetilde h(w)=(\mathrm DG(x))^{-1}.
	\]
	把零测例外集在 \(\mathbb R^n\times\mathbb R^k\) 中作 Fubini 切片。若 \(k>0\)，对几乎所有 \(y\)，在 \(B_y=\{z:(y,z)\in B\}\) 的几乎所有点，映射 \(z\mapsto\widetilde h(y,z)\) 的导数是 \(\mathrm D\widetilde h(y,z)\iota\)。在 \(B_y\) 上它恰为 \(h_y(z)=h(y,z)\)，并将 \(B_y\) 一一映到 \(E\cap f^{-1}(y)\)。由\cref{g3:cor:area-formula-subsets}的单射加权面积公式，
	\begin{equation}\label{g3:eq:coarea-fiber-area}
		\int_{E\cap f^{-1}(y)}g\,\mathrm d\mathcal H^k
		=\int_{B_y}g(h(y,z))J_kh_y(z)\,\mathrm dz.
	\end{equation}
	这里 \(J_kh_y\) 是 \(\sqrt{\det((\mathrm D_zh_y)^{\mathsf T}\mathrm D_zh_y)}\)。面积公式允许定义域为可测子集，且导数在密度点与延拓选择无关。\(k=0\) 时，\(B_y\) 为空集或一个点，\(h_y\) 的面积因子取一，式\eqref{g3:eq:coarea-fiber-area}就是单点的计数恒等式。
	在上述几乎所有点应用式\eqref{g3:eq:coarea-linear-block-identity}，得到
	\[
		J_kh_y(z)=\frac{J_nf(h(y,z))}{|\det\mathrm DG(h(y,z))|}.
	\]
	将这一式代回式\eqref{g3:eq:coarea-fiber-area}，在 \(y\) 上积分，并对方阵双 Lipschitz 映射 \(G\) 换元，便有
	\[
		\begin{aligned}
		\int_{\mathbb R^n}\int_{E\cap f^{-1}(y)}g\,\mathrm d\mathcal H^k\,\mathrm dy
		&=\int_B\frac{g(h(w))J_nf(h(w))}{|\det\mathrm DG(h(w))|}\,\mathrm dw\\
		&=\int_Eg(x)J_nf(x)\,\mathrm dx.
		\end{aligned}
	\]
	这些交换先对 Borel \(g\) 使用非负积分定理；一般 Lebesgue 可测 \(g\) 只与 Borel 代表相差于源零测集，\cref{g3:lem:coarea-covering-inequality}使相应切片差异也几乎处处为零。因此同一等式适用。
\end{proof}

\begin{theorem}[余面积公式]\label{g3:thm:lipschitz-coarea-formula}
	设 \(f:\mathbb R^m\to\mathbb R^n\) 为 Lipschitz 映射，\(m\ge n\ge1\)，\(g:\mathbb R^m\to[0,+\infty]\) 为 Lebesgue 可测函数。则纤维积分在几乎处处意义下可测，并有
	\begin{equation}\label{g3:eq:lipschitz-coarea-formula}
		\int_{\mathbb R^m}g(x)J_nf(x)\,\mathrm dx
		=\int_{\mathbb R^n}\left(\int_{f^{-1}(y)}g(x)\,\mathrm d\mathcal H^{m-n}(x)\right)\mathrm dy.
	\end{equation}
	特别地，对任意 Lebesgue 可测 \(E\subset\mathbb R^m\)，可将左端限制到 \(E\)，并将右端内层限制到 \(E\cap f^{-1}(y)\)。若 \(g\) 为实值或复值，且 \(\int|g|J_nf<\infty\)，则同一公式对绝对可积的切片积分成立。
\end{theorem}
\begin{proof}
	由 Rademacher 定理，不可微点组成一个 \(\mathcal L^m\) 零测集。加上\cref{g3:lem:coarea-regular-pieces}中紧集分片留下的零测集，记为 \(N\)。于是
	\[
		\mathbb R^m=N\cup C\cup\bigcup_jE_j,
	\]
	除去重叠的零测点后，各部分互不相交，其中 \(C\) 是秩亏可微点集合。左端在 \(N\) 上积分为零，在 \(C\) 上因 \(J_nf=0\) 也为零。右端方面，\cref{g3:lem:coarea-covering-inequality,g3:lem:coarea-critical-slices}分别保证 \(N\) 和 \(C\) 与几乎所有纤维的交具有零 \(\mathcal H^k\) 测度。故两端都只须对 \(E_j\) 求和。逐片应用\cref{g3:lem:coarea-on-bilipschitz-piece}，再用非负积分与可数求和交换，即得式\eqref{g3:eq:lipschitz-coarea-formula}。可测性已由\cref{g3:lem:coarea-slice-measurability}建立。
	把 \(g\) 换成 \(g\mathbf1_E\)，得到集合限制版。对一般实值或复值函数先对 \(|g|\) 使用非负公式，得到几乎所有切片的绝对可积性，再对正负部分或实虚部使用同一等式，完成证明。
\end{proof}

\subsection{光滑水平集及非光滑切片}
\label{g3:subsec:2004:04}
\begin{corollary}\label{g3:cor:coarea-almost-every-slice}
	设 \(f:\mathbb R^m\to\mathbb R^n\) 为 Lipschitz 映射，\(m\ge n\ge1\)，\(k=m-n\)。对几乎所有 \(y\)，纤维 \(f^{-1}(y)\) 除去一个 \(\mathcal H^k\) 零测集后，是可数个 \(\mathbb R^k\) 中子集的 Lipschitz 像之并，而且 \(\mathcal H^k\) 在该纤维上局部有限。\(k=0\) 时，这表示纤维在每个有界集内只有有限个点。
\end{corollary}
\begin{proof}
	上述证明中的 \(E_j\cap f^{-1}(y)=h_{j,y}(B_{j,y})\) 已给出这些参数片。把每个定义域再与有界球相交，就可令每片的参数域有界；秩亏点和源零测集的切片余项几乎处处为零。
	对正整数 \(R\)，余面积公式与 \(J_nf\le L^n\) 给出
	\[
		\int_{\mathbb R^n}\mathcal H^k(B(0,R)\cap f^{-1}(y))\,\mathrm dy
		\le L^n\mathcal L^m(B(0,R))<\infty.
	\]
	对可数个 \(R\) 同时去掉零测例外集，就得到局部有限性。\(\mathcal H^0\) 是计数测度，所以 \(k=0\) 的断言也是同一个估计的结果。
\end{proof}

上一推论给出的可数 Lipschitz 参数化，正是\secref{g3:sec:2005}将要研究的可整流结构。它比“水平集是一张光滑子流形”弱，却足以定义和比较几乎处处切片的面积。

\ProseParagraphBreak

若 \(f\) 在开集 \(\Omega\subset\mathbb R^m\) 上为 \(C^1\)，在每个满秩点选择上述补充坐标后，\cref{g3:cor:banach-implicit-parameters}的有限维情形说明局部纤维是 \(k\) 维 \(C^1\) 子流形；其切空间为 \(\ker\mathrm Df\)。因此对一个正则值 \(y\)，整个 \(f^{-1}(y)\cap\Omega\) 是 \(C^1\) 子流形。对一般 \(y\)，令 \(C=\{x\in\Omega:\operatorname{rank}\mathrm Df(x)<n\}\)，则
\[
	f^{-1}(y)\cap\Omega
	=\bigl(f^{-1}(y)\cap(\Omega\setminus C)\bigr)
	\cup\bigl(f^{-1}(y)\cap C\bigr).
\]
第一部分是 \(C^1\) 子流形；第二部分对几乎所有 \(y\) 具有零 \(\mathcal H^k\) 测度。这里把 \(\Omega\) 用可数个相对紧球覆盖，在每个球上 \(\mathrm Df\) 有界，故 \(f\) 为 Lipschitz；逐坐标延拓后应用前述结论，再取可数并即可。若 \(f\) 光滑，第一部分也光滑。这一层面的结论只需 \(C^1\)，不要求每个纤维都避开临界点。

\begin{example}\label{g3:exm:coarea-exceptional-fiber}
	设 \(m\ge2\)，令 \(f(x)=\max\{x_1,0\}\)。它是 \(1\)-Lipschitz，且 \(J_1f=1\) 于 \(x_1>0\)，\(J_1f=0\) 于 \(x_1<0\)。当 \(y>0\) 时，纤维是超平面 \(\{x_1=y\}\)；当 \(y<0\) 时为空；但 \(f^{-1}(0)=\{x_1\le0\}\) 含有开集，因而具有无穷的 \((m-1)\) 维 Hausdorff 测度。对任意非负可测 \(g\)，余面积公式恰为
	\[
		\int_{\{x_1>0\}}g\,\mathrm dx
		=\int_0^\infty\int_{\mathbb R^{m-1}}g(y,z)\,\mathrm dz\,\mathrm dy.
	\]
	异常纤维发生在单个目标值零处，不影响外层 Lebesgue 积分。这个例子显示秩亏区域可以有正的源体积，其像所对应的异常切片仍可在余面积积分中消失。
\end{example}

\begin{example}\label{g3:exm:coarea-polar-integration}
	令 \(m\ge2\)、\(f(x)=|x|\)。在原点以外有 \(|\nabla f|=1\)，故对任意非负可测 \(g\)，
	\[
		\int_{\mathbb R^m}g(x)\,\mathrm dx
		=\int_0^\infty\int_{S^{m-1}(r)}g\,\mathrm d\mathcal H^{m-1}\,\mathrm dr
		=\int_0^\infty r^{m-1}\int_{S^{m-1}}g(r\theta)\,\mathrm d\mathcal H^{m-1}(\theta)\,\mathrm dr.
	\]
	第二个等号使用 Hausdorff 测度在伸缩下乘以 \(r^{m-1}\)。因此极坐标积分可理解成先对距离函数的每个水平球面积分，再对半径积分。
\end{example}

在实际换元中，还常把 Jacobian 移到纤维积分一边。令 \(R_f=\{J_nf>0\}\)，对非负可测 \(u\) 把余面积公式用于 \(g=\mathbf1_{R_f}u/J_nf\)，就得到
\[
	\int_{R_f}u\,\mathrm dx
	=\int_{\mathbb R^n}\int_{R_f\cap f^{-1}(y)}\frac{u}{J_nf}\,\mathrm d\mathcal H^k\,\mathrm dy.
\]
若 \(J_nf>0\) 几乎处处，左端可写成整个源空间的积分；若秩亏集合有正测度，则必须保留这个限制。线性模型和 \(C^1\) 水平集的阅读可参见\cite[第10节，式10.1--10.6，第53--56页]{Simon1983}。
% 纲要条目排为编号小节标题，并进入目录与书签；说明段仍保留于纲要注释中。

% !TeX root = ../../Book3.tex
% 纲要标题：### 20.5 可整流集的切向结构
% 纲要位置：几何学计划纲要.md，第 2955 行。
% 正文状态：初稿；尚未正式审核。

% 纲要内容开始
% 
% * 可数 Lipschitz 参数片
% * 近似切平面的存在唯一性
% * 密度和重数
% * 可整流分解与纯不可整流集
% 
% ---
% 
% 纲要内容结束

\section{可整流集的切向结构}
\label{g3:sec:2005}

面积公式让我们能够按参数计算一个 Lipschitz 片的测度，但一个几何集合可以由许多相互交叠的片组成。本节研究去掉参数之后还留下什么：在几乎每一点，把集合放大，其面积分布趋向一个重数为一的平面。为证明这一点，还要控制其他参数片在该点附近贡献的面积。

\ProseParagraphBreak
以下固定 \(1\le m\le n\)。集合的测度均取 \(\mathcal H^m\)，其中的“几乎处处”也以此测度为准；参数域中的零测集则以 \(\mathcal L^m\) 为准。两者通过面积公式相联系。

\subsection{可数 Lipschitz 参数片}
\label{g3:subsec:2005:01}
\begin{definition}\label{g3:def:countably-rectifiable-set}
	一个 \(\mathcal H^m\)-可测集 \(E\subset\mathbb R^n\) 称为\term[keshu kezhengliu ji]{可数 \(m\)-可整流集}[countably \(m\)-rectifiable set]，若存在可数个 Lipschitz 映射 \(f_i:A_i\subset\mathbb R^m\to\mathbb R^n\)，使
	\[
		\mathcal H^m\left(E\setminus\bigcup_{i=1}^{\infty}f_i(A_i)\right)=0.
	\]
	利用\cref{g3:prop:lipschitz-extension-density}，可以等价地要求 \(A_i=\mathbb R^m\)。本节简称这种集合为可整流集；若需要总测度有限或测度局部有限，将另行写明。
\end{definition}

光滑子流形、Lipschitz 图像、可数条线段的并都是可整流集。可整流性容许尖点、折痕以及片之间的相交；它要求的是由低维参数覆盖，而不是整个集合拥有一个光滑坐标系。

\begin{proposition}[双 Lipschitz 片表示]\label{g3:prop:rectifiable-bilipschitz-cover}
	每个可数 \(m\)-可整流集都有 \(\sigma\)-有限的 \(\mathcal H^m\) 测度。并且除去一个 \(\mathcal H^m\)-零测集后，它可写成两两不交的可数个 Borel 集 \(S_j\) 的并，其中每个 \(S_j\) 都是一个有界 Borel 参数集上的双 Lipschitz 映射的像。
\end{proposition}
\begin{proof}
	先把每个参数域分成有界部分。由\cref{g3:prop:hausdorff-lipschitz,g3:thm:hausdorff-lebesgue}，其像的 \(\mathcal H^m\) 测度有限，故 \(E\) 的测度 \(\sigma\)-有限。

	把参数映射延拓到全空间。不可微点集为 \(\mathcal L^m\)-零测集，其像为 \(\mathcal H^m\)-零测集；可微但秩小于 \(m\) 的点集，其像也由面积公式得到零测度。因此只须保留满秩部分。应用\cref{g3:cor:full-rank-bilipschitz-pieces,g3:cor:compact-linearization-pieces}，它的像除去零测集后由可数个紧的双 Lipschitz 片覆盖。

	对这些紧像逐片与 \(E\) 相交。在每个有限测度片上，\(\mathcal H^m\) 的 Borel 正则性与完备性给出包含于这个交集的 Borel 子集，遗漏的测度为零：具体地，先取交集的 Borel 外包，再在该紧像内对差集取 Borel 外包并作差。有限测度保证两个外包都可取成同测度。于是可用 Borel 子片替换这些交集。按顺序扣除此前的片，使像集两两不交；其在紧参数片中的原像仍为 Borel 集。限制不会破坏双 Lipschitz 性。
\end{proof}

这里的 \(\sigma\)-有限不能换成局部有限。例如，平面内所有水平直线 \(\{y=q\}\)、\(q\in\mathbb Q\) 的并是可数 \(1\)-可整流集，但每个非空开球内的长度都是无穷。另一方面，给一条线段添上平面内一个可数稠密集，所得集合仍可整流，并且长度测度没有改变；它的闭包却是整个平面。这说明测度意义的局部形状不能从集合闭包单独读出。

\ProseParagraphBreak
后面研究的局部有限条件是
\[
	\mathcal H^m(E\cap K)<\infty
	\qquad\text{对每个紧集 }K\subset\mathbb R^n.
\]
它使 \(\mu=\mathcal H^m\llcorner E\) 成为 Radon 测度，其中 \(\llcorner\) 表示测度的限制。对 \(\mathcal H^m\)-可测的 \(E\)，这里取限制测度在 Borel 集上的值；取其完备化不改变下述结论。
\symidx{\(\mathcal H^m\llcorner E\)：Hausdorff 测度在集合 \(E\) 上的限制}

\subsection{近似切平面的存在唯一性}
\label{g3:subsec:2005:02}
\begin{definition}\label{g3:def:approximate-tangent-plane}
	设 \(\mu=\mathcal H^m\llcorner E\) 局部有限，\(x\in\mathbb R^n\)，\(P\subset\mathbb R^n\) 为过原点的 \(m\) 维线性子空间。如果对每个连续紧支撑函数 \(\varphi\in C_c(\mathbb R^n)\)，都有
	\begin{equation}\label{g3:eq:rectifiable-blowup}
		\lim_{r\downarrow0}r^{-m}\int_E
		\varphi\!\left(\frac{y-x}{r}\right)\,\mathrm d\mathcal H^m(y)
		=\int_P\varphi(z)\,\mathrm d\mathcal H^m(z),
	\end{equation}
	就称 \(P\) 为 \(E\) 在 \(x\) 的\term[jinsi qiepingmian]{近似切平面}[approximate tangent plane]，记为 \(\operatorname{Tan}^m(E,x)\)。位于原空间中的仿射切平面是 \(x+P\)。
\end{definition}
\symidx{\(\operatorname{Tan}^m(E,x)\)：以原点为基点的近似切平面}

定义比较的是缩放后的面积分布：左边把 \(x\) 移到原点、把长度放大 \(1/r\)，同时把面积除以 \(r^m\)。采用测试函数，既能记录形状，也能记录单位面积的质量。这个定义与下面的密度、方向锥描述见 \cite[第11节，定义11.2及定理11.6，第59--63页]{Simon1983}。

\begin{proposition}\label{g3:prop:approximate-tangent-consequences}
	近似切平面一旦存在就唯一，并且
	\begin{equation}\label{g3:eq:rectifiable-density-one}
		\lim_{r\downarrow0}
		\frac{\mathcal H^m(E\cap B(x,r))}{\omega_mr^m}=1.
	\end{equation}
	对每个 \(\alpha>0\)，还有
	\begin{equation}\label{g3:eq:rectifiable-cone-concentration}
		\lim_{r\downarrow0}r^{-m}\mathcal H^m
		\{y\in E\cap B(x,r):\operatorname{dist}(y-x,P)\ge\alpha|y-x|\}=0.
	\end{equation}
\end{proposition}
\begin{proof}
	若两个不同的 \(m\) 平面 \(P,Q\) 都满足定义，取 \(z\in P\setminus Q\)，以及支撑在一个与 \(Q\) 不交的小球内、在 \(z\) 附近为正的非负连续函数。它在 \(P\) 上的积分为正，在 \(Q\) 上的积分为零，与同一个极限矛盾。

	用连续函数从内外逼近单位球的示性函数。\(P\) 上的单位球边界是 \(\mathcal H^m\)-零测集，内外球的测度分别趋向 \(\omega_m\)，故得到密度公式。对于锥结论，令
	\[
		F=\{z:|z|\le1,\ \operatorname{dist}(z,P)\ge\alpha|z|\}.
	\]
	这是紧集，且 \(F\cap P\subset\{0\}\)，所以 \(\mathcal H^m(F\cap P)=0\)。取等于一于 \(F\)、支撑在其逐渐缩小邻域内的连续截断函数。先使用测试函数极限，再缩小邻域，得到左边的上极限为零。
\end{proof}

证明存在性时，先研究一个片，然后排除别的片在该点的质量。后一步所需的估计如下。

\begin{lemma}[其他片的质量密度]\label{g3:lem:off-sheet-density-zero}
	设 \(S\) 是有界参数集的双 Lipschitz 像，且为 Borel 集。若 \(\nu\) 为 Radon 测度，\(\nu(S)=0\)，则
	\[
		\lim_{r\downarrow0}\frac{\nu(B(x,r))}{r^m}=0
		\quad\text{对 }\mathcal H^m\text{-几乎每个 }x\in S.
	\]
\end{lemma}
\begin{proof}
	记参数映射及其逆的 Lipschitz 常数为 \(L,L'\)。对任意球 \(B(z,r)\)，若它与 \(S\) 相交，选一个交点作为比较点，其余交点的参数均落在半径 \(2L'r\) 的球内。因此
	\[
		\mathcal H^m(S\cap B(z,r))\le C r^m,
		\qquad C=L^m\omega_m(2L')^m.
	\]
	同一估计也适用于闭球。

	固定紧集 \(K\subset S\)、\(t>0\) 和 \(\varepsilon>0\)。因 \(\nu(K)=0\)，可取开邻域 \(O\supset K\)，使 \(\nu(O)<\varepsilon\)。对于满足
	\(\limsup_{r\downarrow0}\nu(B(x,r))/r^m>t\) 的 \(x\in K\)，有任意小的闭球 \(\overline B(x,r)\subset O\) 满足 \(\nu(\overline B(x,r))>tr^m\)。对这些球应用五倍球覆盖引理，取两两不交的 \(B_i=\overline B(x_i,r_i)\)，其五倍球覆盖全部这样的点。于是该坏集的 \(\mathcal H^m\) 外测度至多为
	\[
		\sum_i C(5r_i)^m
		\le\frac{C5^m}{t}\sum_i\nu(B_i)
		\le\frac{C5^m}{t}\nu(O)
		<\frac{C5^m}{t}\varepsilon.
	\]
	令 \(\varepsilon\downarrow0\)，再取可数个正有理数 \(t\)。最后用紧子集把有限测度的 \(S\) 穷尽至零测集，即得结论。
\end{proof}

\begin{theorem}[可整流集的切平面]\label{g3:thm:rectifiable-tangent-plane}
	设 \(E\subset\mathbb R^n\) 可数 \(m\)-可整流，且 \(\mathcal H^m\llcorner E\) 局部有限。则对 \(\mathcal H^m\)-几乎每个 \(x\in E\)，存在唯一的近似切平面，并且该点的 \(m\) 维密度等于一。

	更具体地，若 \(A\subset\mathbb R^m\) 为 Lebesgue 可测集，\(f:A\to S\subset E\) 是双 Lipschitz 参数片，则对几乎每个 \(a\in A\)，
	\begin{equation}\label{g3:eq:tangent-parametric-image}
		\operatorname{Tan}^m(E,f(a))=\mathrm Df(a)(\mathbb R^m).
	\end{equation}
	这里 \(\mathrm Df\) 由可测参数域上的延拓独立性确定。
\end{theorem}
\begin{proof}
	先只考虑一个有界的双 Lipschitz 片 \(S=f(A)\)。把 \(f\) 延拓到全空间。在几乎每个 \(a\in A\)，\(a\) 是 \(A\) 的密度为一的点，延拓可微，\(T=\mathrm Df(a)\) 的秩为 \(m\)，且
	\(J(z)=J_mf(z)\mathbf1_A(z)\) 在 \(a\) 是 Lebesgue 点，取值 \(J(a)=J_m(T)\)。满秩性来自片上的反向 Lipschitz 估计和密度点处的无大空洞估计。

	固定这样的 \(a\)，令 \(x=f(a)\)。对 \(\varphi\in C_c(\mathbb R^n)\)，单射面积公式和代换 \(u=a+rz\) 给出
	\[
		r^{-m}\int_S\varphi\!\left(\frac{y-x}{r}\right)\,\mathrm d\mathcal H^m(y)
		=\int_{(A-a)/r}
		\varphi\!\left(\frac{f(a+rz)-f(a)}r\right)J_mf(a+rz)\,\mathrm dz.
	\]
	若 \(\operatorname{supp}\varphi\subset B(0,R)\)，片的逆 Lipschitz 常数 \(L'\) 使右端的非零部分包含在固定球 \(|z|\le L'R\) 中。在这个球上，可微性给出
	\[
		\frac{f(a+rz)-f(a)}r\longrightarrow Tz
		\quad\text{一致成立}.
	\]
	而 Lebesgue 点性质给出 \(J(a+r\,\cdot)\to J_m(T)\) 在该球上的 \(L^1\) 收敛。于是右端趋向
	\[
		J_m(T)\int_{\mathbb R^m}\varphi(Tz)\,\mathrm dz
		=\int_{T(\mathbb R^m)}\varphi\,\mathrm d\mathcal H^m.
	\]
	最后一个等式是线性面积公式；积分可写为全空间，因为 \(|Tz|\ge |z|/L'\)。参数中的例外零测集经 Lipschitz 映射仍为 \(\mathcal H^m\)-零测集。

	现在用\cref{g3:prop:rectifiable-bilipschitz-cover} 把 \(E\) 分为两两不交的片 \(S_j\)，余下零测集不贡献测度。对固定的 \(j\)，令
	\(\nu_j=\mathcal H^m\llcorner(E\setminus S_j)\)。它局部有限，且 \(\nu_j(S_j)=0\)。由\cref{g3:lem:off-sheet-density-zero}，在 \(S_j\) 的几乎每一点，
	\[
		r^{-m}\left|\int_{E\setminus S_j}
		\varphi\!\left(\frac{y-x}{r}\right)\,\mathrm d\mathcal H^m(y)\right|
		\le \|\varphi\|_\infty r^{-m}\nu_j(B(x,Rr))
		\longrightarrow0.
	\]
	所以整个 \(E\) 与该片有相同的放大极限。对可数个片取并，得到存在性；唯一性和密度结论由\cref{g3:prop:approximate-tangent-consequences} 给出。对任意另给的双 Lipschitz 子片，同一论证也适用，因而公式不依赖最初选用的分解。
\end{proof}

切平面条件还能够反过来识别可整流性。下面给出方向锥转化为图像的具体步骤。

\begin{theorem}[切平面判别]\label{g3:thm:tangent-rectifiability-criterion}
	设 \(\mathcal H^m\llcorner E\) 局部有限。如果 \(E\) 在 \(\mathcal H^m\)-几乎每一点都存在\cref{g3:def:approximate-tangent-plane} 意义下的近似切平面，则 \(E\) 可数 \(m\)-可整流。
\end{theorem}
\begin{proof}
	在切平面 \(P_x\) 存在的点，由密度和锥集中结论，可选整数 \(k\)，使每个 \(0<r<1/k\) 都满足
	\begin{align*}
		\mathcal H^m(E\cap B(x,r))&\ge(\omega_m/2)r^m,\\
		\mathcal H^m\{z\in E\cap B(x,r):
		\operatorname{dist}(z-x,P_x)\ge |z-x|/4\}&\le\varepsilon_0r^m,
	\end{align*}
	其中固定 \(\varepsilon_0=\omega_m/(4\cdot32^m)\)。

	用正交投影矩阵表示 \(m\) 平面。这些矩阵组成紧集，因为它们恰是满足
	\(\Pi^T=\Pi\)、\(\Pi^2=\Pi\)、\(\operatorname{tr}\Pi=m\) 的矩阵，且范数有界。因此可取有限个平面 \(P_j\)，使每个 \(P_x\) 都满足
	\(\|\Pi_{P_x}-\Pi_{P_j}\|<1/16\)。按 \(k,j\) 分组，再与直径小于 \(1/(2k)\) 的可数个小方块相交。

	设 \(x,y\) 属于同一组，\(d=|x-y|>0\)。若
	\(\operatorname{dist}(y-x,P_j)\ge d/2\)，则
	\(\operatorname{dist}(y-x,P_x)\ge7d/16\)。对每个 \(z\in B(y,d/16)\)，有
	\[
		\operatorname{dist}(z-x,P_x)>6d/16,\qquad
		|z-x|<17d/16.
	\]
	所以该球包含于以 \(x\) 为顶点的上述坏锥区域，且包含于 \(B(x,2d)\)。在 \(y\) 处用下密度界，在 \(x\) 处用锥误差界，得到
	\[
		\frac{\omega_m}{2}\left(\frac d{16}\right)^m
		\le\varepsilon_0(2d)^m
		=\frac{\omega_m}{4}\left(\frac d{16}\right)^m,
	\]
	矛盾。因此同一组内总有
	\[
		|\Pi_{P_j^\perp}(y-x)|<\frac12|y-x|,
		\qquad
		|\Pi_{P_j}(y-x)|>\frac{\sqrt3}{2}|y-x|.
	\]
	投影到 \(P_j\) 在这一组上单射，其逆为 Lipschitz；垂直分量因而是投影像上的 Lipschitz 函数。用分量延拓把它延拓到整个 \(P_j\)，这组点便包含在一个 Lipschitz 图像中。组的总数可数，并覆盖所有具有切平面的点，故得到所需参数覆盖。
\end{proof}

\subsection{密度和重数}
\label{g3:subsec:2005:03}
上一小节的密度一描述的是几何集合本身。若还要记录曲面经过了几次，或在不同位置分配不同质量，就应把重数写进测度。

\begin{theorem}[带重数的放大极限]\label{g3:thm:weighted-rectifiable-density}
	设 \(E\) 可数 \(m\)-可整流，且 \(\mathcal H^m\llcorner E\) 局部有限。设 \(\theta:E\to(0,\infty)\) 可测、几乎处处有限，并且
	\(\theta\in L^1_{\mathrm{loc}}(\mathcal H^m\llcorner E)\)。令
	\(\mu_\theta=\theta\,\mathcal H^m\llcorner E\)。则对 \(\mathcal H^m\)-几乎每个 \(x\in E\)，
	\begin{equation}\label{g3:eq:weighted-rectifiable-blowup}
		\lim_{r\downarrow0}r^{-m}\int_E
		\varphi\!\left(\frac{y-x}{r}\right)\theta(y)\,\mathrm d\mathcal H^m(y)
		=\theta(x)\int_{\operatorname{Tan}^m(E,x)}
		\varphi\,\mathrm d\mathcal H^m
	\end{equation}
	对每个 \(\varphi\in C_c(\mathbb R^n)\) 成立。特别地，
	\begin{equation}\label{g3:eq:weighted-density-value}
		\Theta^m(\mu_\theta,x)
		=\lim_{r\downarrow0}\frac{\mu_\theta(B(x,r))}{\omega_mr^m}
		=\theta(x)
	\end{equation}
	几乎处处成立。
\end{theorem}
\begin{proof}
	在有界双 Lipschitz 片 \(S=f(A)\) 上，带权面积公式把左端化为参数积分，其中的密度变为
	\[
		h(a)=\theta(f(a))J_mf(a)\mathbf1_A(a).
	\]
	该函数可积，因为 \(S\) 有界而 \(\mu_\theta\) 局部有限。在 \(h\) 的 Lebesgue 点，上一小节的放大论证原样给出
	\(\theta(f(a))\mathcal H^m\llcorner\mathrm Df(a)(\mathbb R^m)\) 这一极限。片外的加权测度
	\(\nu=\mu_\theta\llcorner(E\setminus S)\) 是 Radon 测度，且 \(\nu(S)=0\)，由\cref{g3:lem:off-sheet-density-zero}，它在该片的几乎每点贡献零极限。可数分片后即得结论。用连续函数从内外逼近球的示性函数，得到密度值。
\end{proof}

当 \(\theta\) 取正整数值时，称它为整数重数。最简单的例子是同一个平面片被参数映射覆盖两次：像集自身的测度为 \(\mathcal H^m\llcorner E\)，参数面积对应的测度却为 \(2\mathcal H^m\llcorner E\)，其密度为二。面积公式中的计数函数正是这种重数的来源。若不同片在同一几何点汇合，先把各片的重数相加，再计算所得测度。

\ProseParagraphBreak
交叉点可能没有单一切平面。例如，两条不同直线的并在交点的放大极限仍是两条直线，其一维密度为二；在其余点，切线存在且密度为一。交点构成长度零测集，所以与几乎处处的结论相容。这个例子也说明，密度为整数和存在单一切平面是不同的信息。

\subsection{可整流分解与纯不可整流集}
\label{g3:subsec:2005:04}
\begin{definition}\label{g3:def:purely-unrectifiable-set}
	一个 \(\mathcal H^m\)-可测集 \(P\subset\mathbb R^n\) 称为\term[chun bukezhengliu ji]{纯 \(m\)-不可整流集}[purely \(m\)-unrectifiable set]，若对每个 Lipschitz 映射 \(f:\mathbb R^m\to\mathbb R^n\)，都有
	\[
		\mathcal H^m(P\cap f(\mathbb R^m))=0.
	\]
	等价地，\(P\) 与每个可数 \(m\)-可整流集的交集均为 \(\mathcal H^m\)-零测集。
\end{definition}

这比“整个集合不可整流”强得多：一个直线段与一个坏集合的并可以不可整流，但直线段部分仍携带可整流的正长度。纯不可整流性要求这种部分完全不存在。

\begin{theorem}[可整流与纯不可整流分解]\label{g3:thm:rectifiable-pure-decomposition}
	设 \(E\subset\mathbb R^n\) 为 \(\mathcal H^m\)-可测集，且 \(\mathcal H^m(E)\) 为 \(\sigma\)-有限。则可写成不交并
	\[
		E=R\mathbin{\dot\cup}P,
	\]
	其中 \(R\) 可数 \(m\)-可整流，\(P\) 纯 \(m\)-不可整流。两部分在相差 \(\mathcal H^m\)-零测集的意义下唯一。
\end{theorem}
\begin{proof}
	先设 \(\mathcal H^m(E)<\infty\)。令 \(s\) 为 \(E\) 的全部可测可整流子集的测度上确界。选取这样的子集 \(R_j\)，使其测度趋向 \(s\)，并令 \(R=\bigcup_jR_j\)。可数个可整流集的并仍可整流，故 \(\mathcal H^m(R)=s\)。若 \(P=E\setminus R\) 与某个 Lipschitz 像的交集测度为正，则把该交集加入 \(R\) 得到更大的可整流子集，与上确界矛盾。该交集可测，因为全空间的连续像是紧像的可数并。

	一般情形先把 \(E\) 分成两两不交的可测有限测度部分，逐个分解，再分别取并。纯不可整流部分的可数并仍纯不可整流。

	若 \(E=R'\mathbin{\dot\cup}P'\) 是另一个分解，则 \(R\setminus R'\subset R\cap P'\)，右边为零测集；反向差集同理。两部分于是都唯一至零测集。
\end{proof}

一个显式的纯不可整流例子有助于说明，有限的整数维 Hausdorff 测度并不保证存在切平面。

\begin{example}[四角 Cantor 集]\label{g3:exm:four-corner-cantor}
	令
	\[
		C=\left\{\sum_{j=1}^{\infty}3\varepsilon_j4^{-j}:
		\varepsilon_j\in\{0,1\}\right\},\qquad K=C\times C\subset[0,1]^2.
	\]
	每一步只保留正方形四角处、边长为原来四分之一的小正方形。\cref{g3:fig:four-corner-cantor} 画出前两步。这个紧集满足
	\[
		0<\mathcal H^1(K)<\infty,
	\]
	但它是纯 \(1\)-不可整流集。

	先证测度断言。第 \(j\) 层有 \(4^j\) 个边长 \(\ell=4^{-j}\) 的小正方形，它们的直径之和为 \(\sqrt2\)，故 \(\mathcal H^1(K)\le\sqrt2\)。在每个第 \(j\) 层小正方形上赋予质量 \(4^{-j}\)，得到 \(K\) 上的概率测度 \(\lambda\)；具体地，把 \([0,1)\) 上一个点的四进位数字 \(0,1,2,3\) 依次对应四种角位，得到映到 \(K\) 的可测映射，再推前 Lebesgue 测度即可。四进位表示不唯一的可数个点任意选定一种表示；前 \(j\) 位固定的数字柱区间长度为 \(4^{-j}\)，给出所需质量。

	若 \(\ell\le r<4\ell\)，一个半径 \(r\) 的球只与数目有统一上界的第 \(j\) 层小正方形相交。事实上，各方形中心至少相距 \(3\ell\)，围绕这些中心的半径 \(\ell\) 开圆盘不交，并包含在半径 \(r+2\ell\) 的圆盘内；比较面积可取数目上界 \(36\)。因此
	\(\lambda(B(x,r))\le36\ell\le36r\)；对 \(r\ge1\) 也成立。任意直径 \(d>0\) 的集合可包含于一个半径略大于 \(d\) 的球，故其 \(\lambda\)-外测度至多为 \(36d\)。对 \(K\) 的任意覆盖求和，得到
	\(1\le36\sum_i\operatorname{diam}U_i\)，从而 \(\mathcal H^1(K)\ge1/36\)。

	接着证明方向锥集中在 \(K\) 的每一点都失败。固定 \(x\in K\) 和任一直线 \(P\)（过原点）。设第 \(j\) 层包含 \(x\) 的正方形边长为 \(\ell\)。对过 \(x\) 的直线 \(x+P\)，四个顶点中至少有一个到该直线的距离不小于 \(\ell/2\)：正方形沿任一单位法向量的投影宽度至少为 \(\ell\)，两个极端投影不可能都距零小于 \(\ell/2\)。

	取这个顶点处再迭代两次所得的小正方形，其边长为 \(\ell/16\)。其中的 \(K\) 的每一点到 \(x+P\) 的距离至少为
	\[
		\frac{\ell}{2}-\frac{\sqrt2\,\ell}{16}>\frac{\ell}{4},
	\]
	而到 \(x\) 的距离小于 \(2\ell\)。由相似性，该小片的长度测度为 \((\ell/16)\mathcal H^1(K)\)。令 \(r=2\ell\)，于是
	\[
		\frac1r\mathcal H^1
		\{y\in K\cap B(x,r):\operatorname{dist}(y-x,P)\ge |y-x|/8\}
		\ge\frac{\mathcal H^1(K)}{32}.
	\]
	沿 \(j\to\infty\)，这个下界不趋于零。因此 \(K\) 处处没有近似切线。

	最后，若 \(K\) 与某个可整流集相交于正长度的集合，则其中有一个正长度的 Borel 双 Lipschitz 子片 \(S\)。在 \(S\) 的几乎每点，\(S\) 有近似切线；由\cref{g3:lem:off-sheet-density-zero}，测度
	\(\mathcal H^1\llcorner(K\setminus S)\) 在这些点的密度为零。所以 \(K\) 在其中几乎每点也有相同的近似切线，与刚才的结论矛盾。故 \(K\) 纯不可整流。
\end{example}

\begin{figure}[htbp]
	\centering
	% !TeX root = ../Book3.tex
\begin{tikzpicture}[x=0.82cm,y=0.82cm]
	\draw[gray,dashed] (0,0) rectangle (4,4);
	\foreach \x in {0,3} {
		\foreach \y in {0,3} {
			\filldraw[fill=blue!18,draw=blue!70!black] (\x,\y) rectangle ++(1,1);
		}
	}
	\node at (2,-0.5) {第一级：\(4\) 个正方形};
	\begin{scope}[xshift=6cm]
		\draw[gray,dashed] (0,0) rectangle (4,4);
		\foreach \x in {0,3} {
			\foreach \y in {0,3} {
				\draw[gray,dotted] (\x,\y) rectangle ++(1,1);
				\foreach \a in {0,0.75} {
					\foreach \b in {0,0.75} {
						\filldraw[fill=blue!30,draw=blue!70!black]
							({\x+\a},{\y+\b}) rectangle ++(0.25,0.25);
					}
				}
			}
		}
		\node at (2,-0.5) {第二级：\(16\) 个正方形};
	\end{scope}
\end{tikzpicture}

	\caption{四角 Cantor 集的前两步。每层的正方形数量乘以四、边长除以四，所以覆盖直径的总和保持有界；不同角位却在每个尺度上阻止质量集中到一条直线。}
	\label{g3:fig:four-corner-cantor}
\end{figure}

分解定理见 \cite[第13节，定理13.1，第70页]{Simon1983}。它说明弱曲面理论为什么需要指定可整流性：面积有限提供质量控制，可整流性进一步提供几乎处处的切向线性空间；重数则在此基础上记录质量。下一章研究函数变差与周长时，这三项将分别承担不同的作用。
% 纲要条目排为编号小节标题，并进入目录与书签；说明段仍保留于纲要注释中。


\begin{chaptersummary}
	Hausdorff 测度由小直径覆盖定义。标准归一化使它在 \(m\) 维平面上恰等于 Lebesgue 测度；等直径不等式给出一侧比较，Vitali 覆盖给出另一侧。五倍球估计还产生最大函数控制，进而得到 Lebesgue 点、密度点和奇异测度的零密度结论。光滑参数的局部线性比较把这种测度与 Riemann 面积联系起来。

	\ProseParagraphBreak

	绝对连续函数可以由其几乎处处导数积分恢复。对 Lipschitz 映射，先沿可数个方向求导，再用弱导数的线性性与有限方向网，把方向极限提升为全微分。密度点保证可测参数域上的微分不依赖延拓；固定误差和尺度后，可微集分为可数个一致线性化片，满秩片进一步具有反向距离控制。

	\ProseParagraphBreak

	面积公式将源域中的 \(J_mf\) 积分写成像上的重数积分。其证明分别处理降秩像的零面积、近线性单射片的面积比较，以及多片的重数求和。余面积公式则把 \(J_nf\) 积分写成沿 \((m-n)\) 维纤维的积分：覆盖估计先处理源零测集，薄盒估计处理秩亏点，补充坐标后的双 Lipschitz 换元处理满秩部分。

	\ProseParagraphBreak

	可整流性提供可数参数片，局部有限的面积测度使放大极限可以在紧支撑测试函数上比较。在几乎每一点，一个片的微分给出单位密度的近似切平面，其他片的质量具有零密度；正权重则把该极限乘以其几乎处处值。纯不可整流部分与所有 Lipschitz 片都只有零测交集。四角 Cantor 集表明，正且有限的整数维 Hausdorff 测度本身仍不足以产生切平面。
\end{chaptersummary}

\BookExercises

\begin{exercise}\label{g3:ex:holder-hausdorff-dimension}
	设 \(f:(X,d_X)\to(Y,d_Y)\) 满足
	\(d_Y(f(x),f(y))\le Ld_X(x,y)^\alpha\)，其中 \(L>0\)、\(0<\alpha\le1\)。设 \(E\subset X\)。对 \(s>0\)，用本章的归一化证明
	\[
		\mathcal H^{s/\alpha}(f(E))
		\le \frac{\omega_{s/\alpha}}{\omega_s}
		2^{s-s/\alpha}L^{s/\alpha}\mathcal H^s(E),
	\]
	并推出 \(\dim_H f(E)\le\alpha^{-1}\dim_HE\)。
	再证明 \(d^\alpha\) 是一个度量，并计算同一集合在 \(d\) 与 \(d^\alpha\) 下的 Hausdorff 维数关系。
\end{exercise}
\begin{solution}
	若 \(\{U_i\}\) 是 \(E\) 的一个直径不超过 \(\delta\) 的覆盖，则 \(\{f(U_i)\}\) 覆盖 \(f(E)\)，且
	\(\operatorname{diam}f(U_i)\le L(\operatorname{diam}U_i)^\alpha\)。
	令 \(t=s/\alpha\)。逐项计算覆盖代价，
	\[
		\omega_t\left(\frac{\operatorname{diam}f(U_i)}2\right)^t
		\le\frac{\omega_t}{\omega_s}2^{s-t}L^t
		\omega_s\left(\frac{\operatorname{diam}U_i}2\right)^s.
	\]
	取所有覆盖的下确界，再令 \(\delta\downarrow0\)，便得测度不等式。若 \(s>\dim_HE\)，则 \(\mathcal H^s(E)=0\)，故
	\(\mathcal H^{s/\alpha}(f(E))=0\)。令 \(s\downarrow\dim_HE\) 即得维数界；维数无穷时该界自然成立。

	对 \(a,b\ge0\)，有 \((a+b)^\alpha\le a^\alpha+b^\alpha\)。当 \(a>0\) 时，可把它归结为
	\((1+t)^\alpha-1\le t^\alpha\)，其导数比较在 \(t>0\) 上成立；\(a=0\) 时直接成立。结合 \(d\) 的三角不等式，得到 \(d^\alpha\) 的三角不等式，其余度量公理显然满足。

	同一个覆盖集在新度量下的直径是旧直径的 \(\alpha\) 次方。因此对 \(t>0\)，实际上有精确关系
	\[
		\mathcal H^t_{d^\alpha,\delta^\alpha}(E)
		=\frac{\omega_t}{\omega_{\alpha t}}2^{\alpha t-t}
		\mathcal H^{\alpha t}_{d,\delta}(E).
	\]
	令 \(\delta\downarrow0\)，得到
	\(\dim_{H,d^\alpha}E=\alpha^{-1}\dim_{H,d}E\)。
	这说明同一拓扑下的度量尺度改变仍可能改变 Hausdorff 维数。
\end{solution}

\begin{exercise}\label{g3:ex:fat-cantor-density}
	从 \([0,1]\) 开始，在第 \(j\) 步从每个剩余闭区间的中央删去一个长度为 \(4^{-j}\) 的开区间，令 \(F\) 为全部剩余集合的交。证明构造可以一直进行，且 \(F\) 是无内点的闭集，\(\mathcal H^1(F)=1/2\)。证明 \(F\) 的几乎每一点既是拓扑边界点，又是测度密度为一的点，并解释二者为何相容。
\end{exercise}
\begin{solution}
	第 \(j\) 步删去 \(2^{j-1}\) 个开区间，删去的总长度为
	\(2^{j-1}4^{-j}=2^{-j-1}\)。第 \(j\) 步之后总剩余长度为
	\[
		1-\sum_{i=1}^j2^{-i-1}=\frac12+2^{-j-1}.
	\]
	对称性保证剩余的 \(2^j\) 个闭区间等长，每个长度为
	\(2^{-j-1}+2^{-2j-1}\)。在下一步这个长度大于 \(4^{-(j+1)}\)，所以总能继续删去中间开区间。

	这些集合是递减的非空紧集，因而交集 \(F\) 紧。剩余分支区间的最大长度趋于零；若 \(F\) 包含一个非退化区间，该区间在每一级都必须落在单个分支区间中，这与长度趋零矛盾。因此 \(F\) 无内点。由有限测度的上连续性及 \(\mathcal H^1=\mathcal L^1\)，
	\[
		\mathcal H^1(F)=\lim_{j\to\infty}\left(\frac12+2^{-j-1}\right)=\frac12.
	\]
	闭且无内点说明每个 \(F\) 中的点都是其拓扑边界点。Lebesgue 密度定理则说明，对几乎每个 \(x\in F\)，
	\[
		\lim_{r\downarrow0}\frac{\mathcal L^1(F\cap(x-r,x+r))}{2r}=1.
	\]
	边界点要求每个邻域都遇到补集；密度一只要求补集在缩小邻域中的长度比例趋于零。补集可以始终非空而比例越来越小，故两个结论相容。
\end{solution}

\begin{exercise}\label{g3:ex:oscillating-graph-length}
	令 \(\Gamma=\{(x,\sin(1/x)):0<x\le1\}\)。求其闭包，证明
	\(\mathcal H^1(\Gamma)=\infty\)，但 \(\dim_H\overline\Gamma=1\)。
\end{exercise}
\begin{solution}
	若图上点列的横坐标趋于正数，极限仍在图上；若横坐标趋于零，纵坐标的任何极限都在 \([-1,1]\) 内。反之，对每个 \(a\in[-1,1]\)，可选 \(t_j\to\infty\) 使 \(\sin t_j=a\)，于是 \((1/t_j,a)\to(0,a)\)。故
	\[
		\overline\Gamma=\Gamma\cup(\{0\}\times[-1,1]).
	\]
	对 \(\delta>0\)，图在 \([\delta,1]\) 上光滑且单射参数化，长度为
	\[
		\int_\delta^1\sqrt{1+x^{-4}\cos^2(1/x)}\,\mathrm dx
		\ge\int_\delta^1x^{-2}|\cos(1/x)|\,\mathrm dx
		=\int_1^{1/\delta}|\cos t|\,\mathrm dt.
	\]
	最后一个积分随 \(\delta\downarrow0\) 发散，因为每个周期上的积分为一个固定正数。故 \(\mathcal H^1(\Gamma)=\infty\)。

	另一方面，\(\Gamma\) 是图在 \([1/k,1]\) 上各部分的可数并，每一部分都是一个有界区间的 Lipschitz 像。因此对每个 \(s>1\)，这些部分以及竖直线段的 \(\mathcal H^s\) 都为零；可数次可加性给出 \(\mathcal H^s(\overline\Gamma)=0\)。所以维数至多为一。闭包包含一个非退化线段，维数又至少为一，故恰为一。临界维数为一，并不要求一维测度有限。
\end{solution}

% 第20章20.2习题片；由章整合者收入练习区。

\begin{exercise}\label{g3:ex:singular-continuous-not-ac}
	从 \(C_0=[0,1]\) 出发，每次删除现有各闭区间的开中间三分之一，所得集合记为 \(C_k\)，令 \(C=\bigcap_kC_k\)。\(C_k\) 有 \(2^k\) 个基本区间，按从左至右的顺序编号为 \(I_{k,j}\)，\(0\le j<2^k\)。定义连续单调函数 \(c_k:[0,1]\to[0,1]\)：它在 \(I_{k,j}\) 上线性地从 \(j2^{-k}\) 增长到 \((j+1)2^{-k}\)，在各间隙上为常数。
	\begin{enumerate}
		\item 证明 \(c_k\) 一致收敛到连续单调函数 \(c\)，且 \(c(0)=0\)、\(c(1)=1\)。
		\item 证明 \(c'=0\) 几乎处处，但 \(c\) 不绝对连续。给出直接违反绝对连续定义的区间族。
	\end{enumerate}
\end{exercise}
\begin{solution}
	构造保持每个已经删除的开区间上的常值，也保持各级基本区间的端点值。若 \(\ell\ge k\)，则在 \(C_k\) 外有 \(c_\ell=c_k\)，而在每个 \(I_{k,j}\) 内，两者都落在长度为 \(2^{-k}\) 的同一值域区间中。因此
	\[
		\|c_\ell-c_k\|_\infty\le2^{-k}.
	\]
	故 \(c_k\) 一致收敛。连续性、单调性及端点值都传给极限。

	\(C_k\) 的总长度为 \((2/3)^k\)，故 \(\mathcal L^1(C)=0\)。若 \(x\notin C\)，它位于某一级删除的开区间中；\(c\) 在这个开邻域上恒定，所以 \(c'(x)=0\)。因而导数几乎处处为零，却有 \(c(1)-c(0)=1\)，已经与绝对连续函数的积分恢复公式不相容。

	也可直接看定义。在 \(C_k\) 的 \(2^k\) 个基本区间上，区间总长度为 \((2/3)^k\to0\)，而每个区间的端点增量恰为 \(2^{-k}\)。于是
	\[
		\sum_{j=0}^{2^k-1}
		|c(\sup I_{k,j})-c(\inf I_{k,j})|=1.
	\]
	取绝对连续定义中的 \(\varepsilon=1/2\)，无论 \(\delta\) 多小，总能取 \(k\) 使这族区间违反相应要求。
\end{solution}

\begin{exercise}\label{g3:ex:directional-not-total-differential}
	在 \(\mathbb R^2\) 上定义
	\[
		f(x,y)=
		\begin{cases}
			\dfrac{x^3}{x^2+y^2},&(x,y)\ne(0,0),\\[4pt]
			0,&(x,y)=(0,0).
		\end{cases}
	\]
	证明 \(f\) 为全局 Lipschitz 函数，在原点的每个双侧方向导数都存在，但在原点没有全微分。说明这个例子与 Rademacher 定理的证明有什么关系。
\end{exercise}
\begin{solution}
	由 \(|f(x,y)|\le\sqrt{x^2+y^2}\)，函数在原点连续。在原点外记 \(r^2=x^2+y^2\)，则
	\[
		\nabla f=\frac{3x^2}{r^2}(1,0)
		-\frac{2x^3}{r^4}(x,y),
		\qquad |\nabla f|\le3+2=5.
	\]
	若连接两点的线段不经过原点，沿线段积分即可得 Lipschitz 常数五。若线段经过原点，就在原点处分成两段，先在避开原点的部分积分，再由连续性取极限；两段长度之和仍等于原线段长度。因此全局估计仍为 \(|f(p)-f(q)|\le5|p-q|\)。

	对任意非零 \(v=(v_1,v_2)\)，有
	\[
		\frac{f(tv)-f(0)}t=\frac{v_1^3}{v_1^2+v_2^2}
		\qquad(t\ne0).
	\]
	故每个双侧方向导数都存在，但这个方向函数不是线性的：在 \(e_1,e_2\) 处的值分别为一和零，在 \(e_1+e_2\) 处却为 \(1/2\)。若全微分存在，由两个偏导数它只能是 \(A(h_1,h_2)=h_1\)，与第三个方向不符。

	因此方向导数存在性的交集本身还不够。Rademacher 定理还利用分布恒等式，排除方向导数不满足 \(g_v=\sum_i v_i g_{e_i}\) 的零测集；本例的原点正属于必须排除的点。在线性关系已经成立后，有限方向网才能给出全微分。
\end{solution}

\begin{exercise}\label{g3:ex:lipschitz-dense-corners}
	设 \(q_1,q_2,\ldots\) 是有理数的无重复枚举，定义
	\[
		f(x)=\sum_{j=1}^\infty2^{-j}\bigl(|x-q_j|-|q_j|\bigr),
		\qquad x\in\mathbb R.
	\]
	证明该级数在每个有界区间上一致收敛，且 \(f\) 为全局 \(1\)-Lipschitz 函数。求出它在每个无理点的导数，以及在每个 \(q_k\) 的左右导数，并说明其不可微集恰好是稠密零测集 \(\mathbb Q\)。
\end{exercise}
\begin{solution}
	反三角不等式给出
	\[
		\bigl||x-q_j|-|q_j|\bigr|\le|x|.
	\]
	所以在 \([-R,R]\) 上，级数尾项由 \(R\sum_{j>N}2^{-j}\) 一致控制。每一项的 Lipschitz 常数为 \(2^{-j}\)，故求和后
	\[
		|f(x)-f(y)|\le\sum_j2^{-j}|x-y|=|x-y|.
	\]

	若 \(x\notin\mathbb Q\)，每一项都在 \(x\) 可微，其导数为 \(2^{-j}\operatorname{sgn}(x-q_j)\)。为了合法交换极限与级数，注意每一项的差商绝对值都不超过 \(2^{-j}\)，与增量大小无关。因此先截取前 \(N\) 项并对它们取极限，再以 \(\sum_{j>N}2^{-j}\) 控制余下的差商；令 \(N\to\infty\)，即得
	\[
		f'(x)=\sum_{j=1}^\infty2^{-j}\operatorname{sgn}(x-q_j)
		\qquad(x\notin\mathbb Q).
	\]

	固定 \(k\)，把第 \(k\) 项单独拿出。余下级数中的每一项均在 \(q_k\) 可微，同一个尾项估计说明其和在 \(q_k\) 可微，导数为
	\[
		b_k=\sum_{j\ne k}2^{-j}\operatorname{sgn}(q_k-q_j).
	\]
	第 \(k\) 项的左右导数分别为 \(-2^{-k}\) 和 \(2^{-k}\)，所以
	\[
		f'_-(q_k)=b_k-2^{-k},\qquad
		f'_+(q_k)=b_k+2^{-k}.
	\]
	两者相差 \(2^{1-k}>0\)，故每个 \(q_k\) 都不可微。结合无理点上的导数公式，不可微集正是 \(\mathbb Q\)。这说明几乎处处可微并不保证存在一个非空开区间，使函数在其中处处可微。
\end{solution}

\begin{exercise}\label{g3:ex:polar-area-multiplicity}
	设整数 \(k\ge1\)，在 \(E=[0,1]\times[0,2\pi k]\) 上定义 \(F(r,\theta)=(r\cos\theta,r\sin\theta)\)。求 \(J_2F\)、\(N(F,E,y)\) 在单位圆盘中的几乎处处值以及原点处的值，比较参数面积与像面积。对 \(\alpha\in\mathbb R\)，在 \(r>0\) 处取权 \(a(r,\theta)=r^\alpha\)，在 \(r=0\) 处补零，求其带权参数面积，并说明何时有限。
\end{exercise}
\begin{solution}
	\(F\) 在这个紧长方形上为 Lipschitz 映射，其两个偏导数为 \((\cos\theta,\sin\theta)\) 与 \((-r\sin\theta,r\cos\theta)\)，故 \(\det\mathrm DF=r\)，\(J_2F=r\)。若 \(0<|y|<1\) 且 \(y\) 不在正实轴上，令其辐角 \(\vartheta\in(0,2\pi)\)，原像恰为
	\[
		(|y|,\vartheta+2\pi j),\qquad j=0,\ldots,k-1.
	\]
	所以重数为 \(k\)。正实轴的非零点还由两个角端点共同表示，重数为 \(k+1\)；这一线段的二维测度为零。原点的原像是整个 \(\{0\}\times[0,2\pi k]\)，故其重数为无穷。单位圆周上的端点约定同样不影响二维积分。因此
	\[
		\int_EJ_2F\,\mathrm d\mathcal L^2
		=\int_0^{2\pi k}\int_0^1r\,\mathrm dr\mathrm d\theta
		=\pi k,
		\qquad\mathcal H^2(F(E))=\pi.
	\]
	带权公式在几乎每个非零像点给出纤维和 \(k|y|^\alpha\)，而直接在参数域计算得
	\[
		\int_EaJ_2F\,\mathrm d\mathcal L^2
		=2\pi k\int_0^1r^{\alpha+1}\,\mathrm dr
		=\begin{cases}
			\dfrac{2\pi k}{\alpha+2},&\alpha>-2,\\[3pt]
			+\infty,&\alpha\le-2.
		\end{cases}
	\]
	原点处的无穷重数位于目标零测集，不贡献无权积分；这里所选权在该纤维全为零，其带权和也为零。\(\alpha\le-2\) 时的发散来自原点附近一整族正半径圆周，而非原点本身。
\end{solution}

\begin{exercise}\label{g3:ex:piecewise-linear-graph-area}
	在正方形 \(Q=[-1,1]^2\) 上令 \(u(x,y)=\max\{0,x,y\}\)，并令 \(G(x,y)=(x,y,u(x,y))\)。求图面的 \(\mathcal H^2\) 测度。再令 \(S=[0,2]\times[0,1]\)、\(\phi(s,t)=(s-1,2t-1)\)，直接计算 \(G\circ\phi\) 的参数面积，核验它与前一个答案一致，并说明折线上的不可微点为何不造成遗漏。
\end{exercise}
\begin{solution}
	三个函数 \(0,x,y\) 都是 Lipschitz 函数，其最大值仍为 Lipschitz 函数。在 \(Q\) 中除去有限条分界线后，可分成
	\[
		Q_0=\{x<0,y<0\},\qquad
		Q_x=\{x>0,x>y\},\qquad
		Q_y=\{y>0,y>x\}.
	\]
	相应梯度依次为 \((0,0),(1,0),(0,1)\)。\(Q_0\) 面积为 \(1\)，而
	\[
		\mathcal L^2(Q_x)=\int_0^1(x+1)\,\mathrm dx=\frac32,
		\qquad\mathcal L^2(Q_y)=\frac32.
	\]
	由图面积公式，
	\[
		\mathcal H^2(G(Q))=1+\frac32\sqrt2+\frac32\sqrt2=1+3\sqrt2.
	\]
	映射 \(\phi\) 为双射仿射变换，\(\mathrm D\phi=\operatorname{diag}(1,2)\)。在 \(\phi^{-1}(Q_0)\) 上，\(G\circ\phi\) 的两个切向量为 \((1,0,0),(0,2,0)\)，面积因子为 \(2\)；在 \(\phi^{-1}(Q_x)\) 上为 \((1,0,1),(0,2,0)\)，面积因子为 \(2\sqrt2\)；在 \(\phi^{-1}(Q_y)\) 上为 \((1,0,0),(0,2,2)\)，面积因子也为 \(2\sqrt2\)。这三个参数集的面积依次为 \(1/2,3/4,3/4\)，因而参数积分为
	\[
		2\cdot\frac12+2\sqrt2\cdot\frac34+2\sqrt2\cdot\frac34=1+3\sqrt2.
	\]
	分界线在参数域中为二维零测集，其在 Lipschitz 图映射下的像也为 \(\mathcal H^2\) 零测集。因此逐片计算没有遗漏正面积部分；仿射重参数化的行列式恰好补偿了参数集面积的改变。
\end{solution}

\begin{exercise}\label{g3:ex:complex-power-area-pushforward}
	把复平面视为 \(\mathbb R^2\)，设 \(k\ge2\)，\(D=\{z:|z|<1\}\)，\(F(z)=z^k\)。证明其参数面积为 \(k\pi\)，并求有限测度
	\[
		\nu(B)\coloneqq\mathcal L^2(\{z\in D:z^k\in B\}),\qquad B\subset\mathbb C\text{ 为 Borel 集},
	\]
	相对于 \(\mathcal L^2\) 的密度。核验该密度在原点附近虽无界但可积，并计算 \(\nu(\{|w|<r\})\)、\(0<r<1\)。
\end{exercise}
\begin{solution}
	\(F\) 在闭单位圆盘上 Lipschitz，其实微分是复数乘法 \(v\mapsto kz^{k-1}v\)。若 \(z\ne0\)，两个奇异值都等于 \(k|z|^{k-1}\)，故
	\[
		J_2F(z)=k^2|z|^{2k-2}.
	\]
	原点的微分为零。每个 \(0<|w|<1\) 恰有 \(k\) 个不同的 \(k\) 次根，所以面积公式给出 \(\int_DJ_2F\,\mathrm d\mathcal L^2=k\pi\)。也可直接积分 \(2\pi k^2\int_0^1r^{2k-1}\,\mathrm dr=k\pi\)。

	给定非负 Borel 函数 \(h\)，在 \(z\ne0\) 处对面积公式取权
	\[
		a(z)=\frac{h(z^k)}{k^2|z|^{2k-2}},
	\]
	并令 \(a(0)=0\)。每个非零 \(w\) 的全部原像有相同模长 \(|w|^{1/k}\)，因此
	\[
		\int_Dh(z^k)\,\mathrm d\mathcal L^2(z)
		=\int_Dh(w)\,\frac1k|w|^{2/k-2}\,\mathrm d\mathcal L^2(w).
	\]
	故所求密度可取
	\[
		\frac{\mathrm d\nu}{\mathrm d\mathcal L^2}(w)
		=\frac1k|w|^{2/k-2}\mathbf1_{\{0<|w|<1\}}.
	\]
	它在原点的取值任意。指数 \(2/k-2<0\) 说明密度无界，但极坐标积分中的径向幂为 \(2/k-1>-1\)，所以可积。更具体地，
	\[
		\nu(\{|w|<r\})
		=\frac{2\pi}{k}\int_0^r s^{2/k-1}\,\mathrm ds
		=\pi r^{2/k}.
	\]
	这也等于原像圆盘 \(\{|z|<r^{1/k}\}\) 的面积。令 \(r\uparrow1\)，总质量为 \(\pi\)，与源圆盘的面积一致；它与 \(k\pi\) 的参数面积分别对应权 \(1\) 和权 \(J_2F\)，含义不同。
\end{solution}

% 第20章第4节习题，由章入口统一整合。

\begin{exercise}\label{g3:ex:coarea-oblique-projection}
	令 \(f:\mathbb R^3\to\mathbb R^2\)、\(f(x_1,x_2,x_3)=(x_1+x_2,x_3)\)。计算 \(J_2f\)，用一个等距参数写出每根纤维，并对单位球 \(B\subset\mathbb R^3\) 直接验证
	\[
		\int_{\mathbb R^2}\mathcal H^1(B\cap f^{-1}(u,v))\,\mathrm du\,\mathrm dv
		=\int_BJ_2f\,\mathrm dx.
	\]
	说明若把纤维的某个非等距参数直接当作弧长，会漏掉哪个因子。
\end{exercise}
\begin{solution}
	有
	\[
		\mathrm Df=\begin{pmatrix}1&1&0\\0&0&1\end{pmatrix},
		\qquad \mathrm Df(\mathrm Df)^{\mathsf T}=\begin{pmatrix}2&0\\0&1\end{pmatrix},
	\]
	所以 \(J_2f=\sqrt2\)。纤维的等距参数化为
	\[
		\gamma_{u,v}(s)=\left(\frac u2+\frac s{\sqrt2},
		\frac u2-\frac s{\sqrt2},v\right),\qquad s\in\mathbb R.
	\]
	确有 \(|\gamma'|=1\)，且 \(|\gamma_{u,v}(s)|^2=u^2/2+v^2+s^2\)。因此
	\[
		\mathcal H^1(B\cap f^{-1}(u,v))
		=\begin{cases}
		2\sqrt{1-u^2/2-v^2},&u^2/2+v^2<1,\\
		0,&u^2/2+v^2\ge1.
		\end{cases}
	\]
	作 \(a=u/\sqrt2\)、\(b=v\)，再在 \((a,b)\) 平面使用极坐标，便得左端为
	\[
		\sqrt2\int_{a^2+b^2<1}2\sqrt{1-a^2-b^2}\,\mathrm da\,\mathrm db
		=4\pi\sqrt2\int_0^1\rho\sqrt{1-\rho^2}\,\mathrm d\rho
		=\frac{4\pi\sqrt2}{3}.
	\]
	右端则是 \(\sqrt2\mathcal L^3(B)=4\pi\sqrt2/3\)，相符。若改用 \(t=x_1-x_2\)，纤维参数成为 \(((u+t)/2,(u-t)/2,v)\)，其速度为 \(1/\sqrt2\)。所以弧长元是 \(\mathrm dt/\sqrt2\)，而不是 \(\mathrm dt\)。
\end{solution}

\begin{exercise}\label{g3:ex:coarea-distance-neighborhood}
	设 \(m\ge2\)，\(K\subset\mathbb R^m\) 非空且紧，令 \(d_K(x)=\operatorname{dist}(x,K)\)。证明 \(d_K\) 是 \(1\)-Lipschitz，且在 \(\mathbb R^m\setminus K\) 的几乎所有点有 \(|\nabla d_K|=1\)。定义
	\[
		V(0)=\mathcal L^m(K),\qquad
		V(r)=\mathcal L^m\{x:d_K(x)<r\}\quad(r>0).
	\]
	证明对每个 \(r>0\)，
	\[
		V(r)-V(0)=\int_0^r\mathcal H^{m-1}\{x:d_K(x)=t\}\,\mathrm dt.
	\]
	由此说明 \(V\) 在每个有界半径区间上绝对连续，并确定其几乎处处导数。
\end{exercise}
\begin{solution}
	对任意 \(p\in K\)，三角不等式给出 \(|x-p|\le|x-y|+|y-p|\)。对 \(p\) 取下确界并交换 \(x,y\)，得到 \(|d_K(x)-d_K(y)|\le|x-y|\)。因此由 Rademacher 定理，\(d_K\) 几乎处处可微，且在可微点有 \(|\nabla d_K|\le1\)。
	设 \(x\notin K\) 是可微点，记 \(d=d_K(x)>0\)。紧性给出 \(p\in K\) 使 \(|x-p|=d\)，令 \(v=(x-p)/d\)。对 \(0<h<d\)，一方面以 \(p\) 作比较得到 \(d_K(x-hv)\le d-h\)，另一方面由 Lipschitz 性得到 \(d_K(x-hv)\ge d-h\)。故
	\[
		d_K(x-hv)=d-h,
		\qquad\nabla d_K(x)\cdot v=1.
	\]
	结合 \(|v|=1\) 和上界，即有 \(|\nabla d_K(x)|=1\)。
	现在把余面积公式用于 \(f=d_K\) 和 \(g=\mathbf1_{\{0<d_K<r\}}\)。源中的被积函数几乎处处等于这个示性函数；在目标值 \(0<t<r\) 处，纤维上的示性函数恒为一，在其他目标值处为零。因此
	\[
		\mathcal L^m\{0<d_K<r\}
		=\int_0^r\mathcal H^{m-1}\{d_K=t\}\,\mathrm dt.
	\]
	左端就是 \(V(r)-V(0)\)。由于 \(K\) 紧，固定 \(R>0\) 时 \(\{d_K<R\}\) 有界，故右端在 \((0,R)\) 上可积。一个可积函数的不定积分在 \([0,R]\) 上绝对连续，于是
	\[
		V'(r)=\mathcal H^{m-1}\{d_K=r\}
	\]
	对几乎所有 \(r>0\) 成立。这里不要求每个距离水平集都光滑；可测性和几乎处处的有限面积由余面积公式提供。
\end{solution}

\begin{exercise}\label{g3:ex:coarea-square-pushforward}
	设 \(m\ge2\)、\(Q=(0,1)^{m-1}\)、\(E=(-1,1)\times Q\)，并令 \(f(s,z)=s^2\)。对任意非负 Lebesgue 可测函数 \(u:E\to[0,\infty]\)，证明
	\[
		\int_Eu(s,z)\,\mathrm ds\,\mathrm dz
		=\int_0^1\frac1{2\sqrt t}
		\left(\int_Q\bigl[u(\sqrt t,z)+u(-\sqrt t,z)\bigr]\,\mathrm dz\right)\mathrm dt.
	\]
	解释 \(t=0\) 处的因子如何处理，并计算 \(E\) 上归一化 Lebesgue 测度经 \(f\) 推前后的密度。
\end{exercise}
\begin{solution}
	可把 \(f\) 延拓为全空间 Lipschitz 函数 \(\widetilde f(s,z)=\min\{s^2,1\}\)，其 Lipschitz 常数至多为二。在 \(E\) 内，\(J_1\widetilde f=2|s|\)，除去超平面 \(s=0\) 后严格为正；这个超平面的 \(\mathcal L^m\) 测度为零。
	将 \(u\) 在 \(E\) 外补零，并在 \(E\cap\{s\ne0\}\) 上取权 \(g=u/(2|s|)\)，其余地方取零。则 \(gJ_1\widetilde f=u\) 几乎处处于 \(E\)。对 \(0<t<1\)，与权函数支集相交的纤维由两片平行平面 \(\{\sqrt t\}\times Q\) 和 \(\{-\sqrt t\}\times Q\) 组成，二者的 \(\mathcal H^{m-1}\) 测度都由 \(\mathrm dz\) 给出。在这两片上，分母均为 \(2\sqrt t\)。代入余面积公式即得所需等式。
	等式的外层积分按 Lebesgue 意义进行，单个值 \(t=0\) 上如何补定义不影响结果。因子 \(t^{-1/2}\) 在 \((0,1)\) 上可积，反映的是靠近临界值时横向变化变慢，并不造成总质量无穷。
	因为 \(\mathcal L^m(E)=2\)，归一化源测度为 \(\tfrac12\mathbf1_E\mathcal L^m\)。对非负可测测试函数 \(\varphi\)，取 \(u(s,z)=\tfrac12\varphi(s^2)\)，上式给出
	\[
		\frac12\int_E\varphi(f(s,z))\,\mathrm ds\,\mathrm dz
		=\int_0^1\varphi(t)\frac1{2\sqrt t}\,\mathrm dt.
	\]
	故推前测度的密度是 \(\mathbf1_{(0,1)}(t)/(2\sqrt t)\)，其积分为一，且在零处没有原子。
\end{solution}

\begin{exercise}\label{g3:ex:circles-zero-density-point}
	在平面中令 \(S_j\) 为圆心 \((2^{-j},0)\)、半径 \(2^{-3j-3}\) 的圆周，\(j\ge1\)，并令
	\(E=\{0\}\cup\bigcup_{j\ge1}S_j\)。
	证明 \(E\) 紧、可数 \(1\)-可整流且长度有限。证明原点属于长度测度的支撑，却有 \(\Theta^1(\mathcal H^1\llcorner E,0)=0\)，因而在原点没有本章定义的近似切线。
\end{exercise}
\begin{solution}
	各圆周紧，其圆心与半径均趋于零。任意取自这些圆周的收敛点列，如果圆周指标有有界子列，极限属于其中某个圆周；否则极限为原点。所以加入原点后的 \(E\) 闭且有界，因而紧。各圆周由光滑参数覆盖，原点为长度零测集，故 \(E\) 可整流。

	这些圆周两两不交。例如相邻圆心的距离为 \(2^{-j-1}\)，两个半径的和小于 \(2^{-j-1}\)；更远圆心的距离更大而后一个半径更小。因此
	\[
		\mathcal H^1(E)=2\pi\sum_{j=1}^{\infty}2^{-3j-3}<\infty.
	\]
	任意原点邻域都包含指标充分大的整个圆周，故原点属于该测度的支撑。

	写 \(c_j=2^{-j}\)、\(\rho_j=2^{-3j-3}\)。因 \(\rho_j\le c_j/4\)，若 \(S_j\cap B(0,r)\ne\varnothing\)，必有 \(c_j-\rho_j<r\)，从而 \(c_j<4r/3\)。于是
	\[
		\mathcal H^1(E\cap B(0,r))
		\le2\pi\sum_{\{j:2^{-j}<4r/3\}}2^{-3j-3}
		\le Cr^3.
	\]
	最后一步是公比 \(1/8\) 的几何级数估计，常数与小半径 \(r\) 无关。除以 \(\omega_1r=2r\) 后趋于零。近似切线要求密度为一，所以在原点不存在。可整流集的切平面定理只在长度几乎处处断言存在性，而原点的长度为零。
\end{solution}

\begin{exercise}\label{g3:ex:crossing-weighted-lines}
	设 \(L_1,L_2\subset\mathbb R^2\) 为两条不同的过原点直线，\(a,b>0\)，令
	\(\mu=a\mathcal H^1\llcorner L_1+b\mathcal H^1\llcorner L_2\)。
	求 \(\Theta^1(\mu,x)\) 在两条直线上的值，以及在原点的所有尺度放大结果。证明即使允许任意正重数，原点的放大结果也不能写成单一直线上的常数倍长度测度。
\end{exercise}
\begin{solution}
	若 \(x\in L_1\setminus\{0\}\)，则它到 \(L_2\) 的距离为正，小球只遇到 \(L_1\)，因而
	\(\mu(B(x,r))=2ar\)，密度为 \(a\)。同理在 \(L_2\setminus\{0\}\) 上密度为 \(b\)。原点处则
	\(\mu(B(0,r))=2(a+b)r\)，密度为 \(a+b\)。

	直线在以原点为中心的伸缩下不变，因此对任意 \(r>0\) 和 \(\varphi\in C_c(\mathbb R^2)\)，
	\[
		r^{-1}\int\varphi(y/r)\,\mathrm d\mu(y)
		=a\int_{L_1}\varphi\,\mathrm d\mathcal H^1
		+b\int_{L_2}\varphi\,\mathrm d\mathcal H^1.
	\]
	所以每个尺度的放大测度都恰为 \(\mu\) 本身。如果它等于 \(c\mathcal H^1\llcorner L\)，则 \(L\) 必须同时包含 \(L_1,L_2\)：否则在某条直线上选择一个不属于 \(L\) 的点，用支撑在与 \(L\) 不交的小球内的非负测试函数，就得到左侧正、右侧零。两条不同直线不能同时包含于一条直线，故这样的表示不存在。密度只给出总质量的尺度系数，不能替代切向方向的信息。
\end{solution}

\begin{exercise}\label{g3:ex:rectifiable-product}
	设 \(E\subset\mathbb R^n\)、\(F\subset\mathbb R^q\) 分别是可数个有界 Borel 双 Lipschitz \(m\) 维片和 \(k\) 维片的并，其中 \(m,k\ge1\)，且相应 Hausdorff 测度局部有限。证明 \(E\times F\) 可数 \((m+k)\)-可整流，并且作为 \(\mathbb R^{n+q}\) 上的 Borel 测度，
	\[
		\mathcal H^{m+k}\llcorner(E\times F)
		=(\mathcal H^m\llcorner E)\otimes(\mathcal H^k\llcorner F).
	\]
	证明在乘积测度几乎处处的 \((x,y)\) 上，其近似切平面为
	\(\operatorname{Tan}^m(E,x)\times\operatorname{Tan}^k(F,y)\)。说明为何对一般可整流集合，应先取双 Lipschitz 分片给出的满测度代表来使用本题。
\end{exercise}
\begin{solution}
	先处理两个双 Lipschitz 片 \(S=f(A)\subset E\)、\(T=g(B)\subset F\)。乘积映射
	\(h(a,b)=(f(a),g(b))\) 仍为双 Lipschitz 映射。它的微分和 Gram 矩阵均按两个因子分块，因此
	\[
		J_{m+k}h(a,b)=J_mf(a)J_kg(b)
	\]
	对几乎每个 \((a,b)\in A\times B\) 成立。对任意非负 Borel 函数 \(\psi\)，带权面积公式、Fubini 定理和两个因子的单射面积公式给出
	\[
		\int_{S\times T}\psi\,\mathrm d\mathcal H^{m+k}
		=\int_S\int_T\psi(x,y)\,\mathrm d\mathcal H^k(y)\,\mathrm d\mathcal H^m(x).
	\]

	按次序从各片扣除之前的片，便可把 \(E,F\) 写成两两不交的 Borel 双 Lipschitz 片的并。相应的乘积片也不交，且覆盖整个 \(E\times F\)。对它们使用上式，再由非负函数的单调收敛交换求和与积分，就得到所需测度恒等式。同时这些乘积参数映射构成可数 Lipschitz 覆盖，所以乘积可整流。

	对 \(\mathbb R^{n+q}\) 中任意紧集，其在两个因子中的投影都紧。所得乘积测度在它上面的质量不超过两个投影测度的乘积，因而局部有限。于是可以应用可整流切平面定理。在几乎处处的参数点，\(h\) 的微分是 \(\mathrm Df\) 与 \(\mathrm Dg\) 的分块直和，其像为
	\[
		\mathrm Df(a)(\mathbb R^m)\times\mathrm Dg(b)(\mathbb R^k).
	\]
	对因子和乘积各使用\cref{g3:thm:rectifiable-tangent-plane}，并用 Fubini 定理排除两个因子的零测例外集，得到所述切平面公式。

	一般可整流集合由\cref{g3:prop:rectifiable-bilipschitz-cover} 给出这样的满测度代表 \(E_0,F_0\)。本题直接适用于 \(E_0\times F_0\)。原定义允许在 \(E,F\) 上添加任意相应维数的零测集，而题中的参数片论证不覆盖这些额外集合的乘积；因此使用本题时应保留明确的代表条件，不能仅凭因子中的零测性就删去任意乘积余集。
\end{solution}
